\chapter{Eta forms and the index of product-type families}\label{ch:eta}
\providecommand{\Ch}{\operatorname{Ch}}
\providecommand{\Chodd}{\operatorname{Ch}_{\mathrm{odd}}}
\providecommand{\Td}{\operatorname{Td}}
\providecommand{\Trbar}{\overline{\operatorname{Tr}}}
\providecommand{\TrR}{\operatorname{Tr}_{\mathrm{R}}}
\providecommand{\res}{\operatorname{res}}
\providecommand{\FP}{\operatorname{FP}}
\providecommand{\wind}{\operatorname{wind}}
\providecommand{\Hardy}{\mathcal{H}}
\providecommand{\GT}{G_{\mathcal{T}}}
\providecommand{\PsiT}[1]{\Psi^{#1}_{\mathcal{T}}}
\providecommand{\Psipt}[1]{\Psi^{#1}_{\mathrm{pt}}}
\providecommand{\PsiptT}[1]{\Psi^{#1}_{\mathrm{pt},\mathcal{T}}}

The index formula for product-type operators involves, besides the familiar
Atiyah--Singer integrand, a `regularized odd Chern character' of the base family,
which we call an eta form. In this chapter we first construct the necessary
regularization of the trace on a compact manifold, using holomorphic families of
pseudodifferential operators: Seeley's theorem on the meromorphic continuation of
traces, the residue trace, the regularized trace and the trace-defect formula (these
are the analogues for compact manifolds of the constructions for the isotropic
algebra in \cref{ch:determinants}). We then specialize to the circle and define eta
forms on the group of invertible Toeplitz operators, compute their exterior
derivative and relate them to the odd Chern forms and the Fredholm determinant of
\cref{ch:index}. The main result, \cref{eta:thm:index}, is an index formula for
fully elliptic operators of Toeplitz product type on $M\times\Sphere$; the proof
uses the product-type calculus of \cref{ch:product}, the weak contractibility of the
Toeplitz group from \cref{ch:bott} and the Atiyah--Singer theorem of
\cref{ch:families}. A version for the full (non-Toeplitz) product-type calculus
follows. The last section discusses, more briefly and partly conjecturally, the
general product-type index problem, a corresponding cohomology theory and the
determinant line bundle. The eta forms here are close relatives of the eta
invariant of Atiyah--Patodi--Singer \cite{Melrose-APS} and of the eta forms of
Bismut and Cheeger \cite{Bismut-Cheeger}; see also \cite{Melrose-Rochon}.

Throughout, $Z$ denotes a compact manifold without boundary of dimension $d\ge1$
and $E\longrightarrow Z$ a complex vector bundle. We fix a smooth positive density
on $Z$ and a Hermitian structure on $E$, so that operators act on $L^2(Z;E)$; the
trace does not depend on these choices.

\section{The trace on operators of low order}\label{eta:sec:trace}

\subsection{Operators of complex order} Recall (\cref{ch:distributions,ch:conormal})
that the classical symbols of order $m$ on $\R^d$ are exactly the functions of the
form $\rho^{-m}b$ with $b\in\CI(\RC{\R^d})$, where $\RC{\R^d}$ is the radial
compactification and $\rho$ is a boundary defining function, for instance
$\rho=\ang{\xi}^{-1}$. This description makes sense for any $m\in\C$: we say that
$a\in\CI(U\times\R^d;M(N,\C))$ is a \emph{classical symbol of order $m\in\C$} if
$a=\rho^{-m}b$ with $b\in\CI(U\times\RC{\R^d};M(N,\C))$. Such a symbol has an
asymptotic expansion $a\sim\sum_{j\ge0}a_{m-j}$ with $a_{m-j}$ homogeneous of degree
$m-j$ in $\abs\xi\ge1$, obtained from the Taylor series of $b$ at the boundary.
Correspondingly $\Psido{m}(Z;E)$, $m\in\C$, denotes the operators whose kernels are
smooth away from the diagonal and which in any coordinate chart $U$ over which $E$ is
trivial, after cut-off, are of the form $\Op(a)$ with $a$ a classical symbol of
order $m$ (left quantization), modulo smoothing operators. The composition theorem
and the symbol calculus of \cref{ch:manifolds} hold without change for complex
orders: $\Psido{m}\circ\Psido{m'}\subset\Psido{m+m'}$, principal symbols multiply and
$\Psido{m}\subset\Psido{m'}$ if $\Re m'-\Re m\in\N_0$ and $m'-m\in\N_0$. We write
\[
\Psido{\Z}(Z;E)=\bigcup_{m\in\Z}\Psido{m}(Z;E)
\]
for the algebra of classical operators of integral order.

\begin{lemma}\label{eta:lem:trace-class}
If $A\in\Psido{m}(Z;E)$ with $\Re m<-d$ then $A$ is of trace class on $L^2(Z;E)$, its
kernel $K_A$ is continuous and
\[
\Tr(A)=\int_Z\tr_E K_A(z,z).
\]
If $U$ is a coordinate chart over which $E$ is trivialized,
$\phi,\psi\in\CIc(U)$ and $A=\phi\Op(a)\psi$ with $a$ a symbol of order $m$, then
\begin{equation}\label{eta:eq:local-trace}
\Tr(A)=(2\pi)^{-d}\int_U\int_{\R^d}\phi(z)\psi(z)\tr a(z,\xi)\,d\xi\,dz.
\end{equation}
\end{lemma}

\begin{proof} Using a partition of unity $\{\phi_i\}$ subordinate to coordinate
charts and $\psi_i\in\CIc$ equal to $1$ near $\supp\phi_i$, write
$A=\sum_i\phi_iA\psi_i+R$ with $R=\sum_i\phi_iA(1-\psi_i)$ smoothing. Smoothing
operators on a compact manifold are of trace class with trace given by the integral
of the kernel over the diagonal, so it suffices to consider $A=\phi\Op(a)\psi$ in a
chart. Write $a=a'\ang\xi^{m/2}$ with $a'=a\ang\xi^{-m/2}$ of order $m/2$. Since the
second factor is a Fourier multiplier, $\Op(a')\Op(\ang\xi^{m/2})=\Op(a)$ exactly in the
left quantization, so $A=T_1T_2$ with $T_1=\phi\Op(a')$, $T_2=\Op(\ang\xi^{m/2})\psi$. The
kernel of $T_2$ is $k(z-z')\psi(z')$ with $\hat k=\ang\xi^{m/2}\in L^2(\R^d)$ since
$\Re m<-d$; hence $\int\int\abs{k(z-z')\psi(z')}^2dz\,dz'=\norm k^2\norm\psi^2$ and
$T_2$ is Hilbert--Schmidt. The kernel of $T_1$ is
$\phi(z)\int e^{i(z-z')\cdot\xi}a'(z,\xi)\dbar\xi$ which, for each $z$, is in $L^2$ in
$z'$ with norm $\norm{a'(z,\cdot)}_{L^2}$, uniformly bounded; as $\phi$ has compact
support $T_1$ is Hilbert--Schmidt. Thus $A$ is of trace class and
$\Tr(A)=\int\int K_{T_1}(z,z')K_{T_2}(z',z)\,dz'\,dz$. Since $a$ is integrable in
$\xi$, locally uniformly in $z$, the kernel
$K_A(z,z')=\phi(z)\int e^{i(z-z')\cdot\xi}a(z,\xi)\dbar\xi\,\psi(z')$ is continuous and
the double integral equals $\int K_A(z,z)\,dz$, which is \eqref{eta:eq:local-trace}.
\end{proof}

\begin{lemma}\label{eta:lem:trace-property}
If $A\in\Psido{m_1}(Z;E)$, $B\in\Psido{m_2}(Z;E)$ and $\Re(m_1+m_2)<-d$ then
$\Tr(AB)=\Tr(BA)$.
\end{lemma}

\begin{proof} If $A$ is smoothing then $AB$ and $BA$ are smoothing and both traces are
the pairing of $K_A(z,z')$ with $K_B(z',z)$, so they agree. In general choose
$\delta>0$ with $\Re(m_1+m_2)+\delta<-d$ and smoothing operators $A_\epsilon\to A$ in
$\Psido{m_1+\delta}$ as $\epsilon\to0$ (for instance by cutting off the local symbols
with $\chi(\epsilon\xi)$, $\chi\in\CIc(\R^d)$ equal to $1$ near $0$). Composition is
continuous and, by \eqref{eta:eq:local-trace}, the trace is continuous on
$\Psido{m}$ for $\Re m<-d$, so $\Tr(AB)=\lim\Tr(A_\epsilon B)=\lim\Tr(BA_\epsilon)=\Tr(BA)$.
\end{proof}

\section{Holomorphic families}\label{eta:sec:holomorphic}

The space $\Psido{m}(Z;E)$ of fixed order has a natural Fr\'echet topology and there is
no difficulty in speaking of holomorphic maps $\Omega\longrightarrow\Psido{m}(Z;E)$ for
$\Omega\subset\C$ open. We need families in which the order varies holomorphically
as well; these are not maps into a fixed space.

\begin{definition}\label{eta:def:hol-family}
Let $\Omega\subset\C$ be open and $\mu:\Omega\longrightarrow\C$ holomorphic. A family
$A(s)\in\Psido{\mu(s)}(Z;E)$, $s\in\Omega$, is a \emph{holomorphic family of order
$\mu$} if
\begin{enumerate}[label=(\roman*)]
\item for every $\chi\in\CI(Z^2)$ vanishing near the diagonal, $\chi K_{A(s)}$ is a
holomorphic map from $\Omega$ to $\CI(Z^2;\Hom(E)\otimes\Omega_R)$, and
\item for every coordinate chart $U$ over which $E$ is trivial and all
$\phi,\psi\in\CIc(U)$,
\[
\phi A(s)\psi=\Op\bigl(\rho^{-\mu(s)}b(s)\bigr)+R(s),
\]
where $b\in\CI(\Omega\times U\times\RC{\R^d};M(N,\C))$ is holomorphic in $s$ with
support in a fixed compact subset of $U$ in $z$, and $R(s)$ is a holomorphic family
of smoothing operators.
\end{enumerate}
\end{definition}

Changing the boundary defining function $\rho$ to $\rho'=f\rho$, $f>0$ smooth, replaces
$b(s)$ by $f^{\mu(s)}b(s)$, which is again smooth and holomorphic in $s$; so the
definition does not depend on $\rho$. Its independence of the coordinates follows,
exactly as for a single operator in \cref{ch:manifolds}, from the fact that the
coordinate-change formula for symbols is an asymptotic sum whose terms are
obtained from the symbol by applying differential operators in $(z,\xi)$ and
multiplying by smooth functions; these operations preserve
$\rho^{-\mu(s)}\CI(\Omega\times U\times\RC{\R^d})$, with holomorphic dependence, and the
asymptotic summation can be carried out holomorphically (Borel's lemma with
cut-offs chosen independently of $s$, uniformly on compact subsets of $\Omega$). For
$\mu$ constant this is the usual notion of a holomorphic family. In the same way one
obtains:

\begin{lemma}\label{eta:lem:hol-composition}
If $A(s)\in\Psido{\mu(s)}(Z;F,G)$ and $B(s)\in\Psido{\nu(s)}(Z;E,F)$ are holomorphic
families on $\Omega$ then $A(s)B(s)$ is a holomorphic family of order $\mu+\nu$.
\end{lemma}

\begin{proof}[Sketch] Off the diagonal this follows from the corresponding property of
compositions with smoothing operators. Near the diagonal, in local coordinates, the
symbol of the composite has the expansion $\sum_\alpha\frac1{\alpha!}\pa_\xi^\alpha a(s)
D_z^\alpha b(s)$, each term lying in $\rho^{-\mu(s)-\nu(s)}\CI$ holomorphically, and the
remainder after $J$ terms satisfies symbol estimates of order $\Re(\mu+\nu)-J$
locally uniformly in $s$; by Cauchy's formula the same holds for its $s$-derivatives,
which gives holomorphy of the composite symbol.
\end{proof}

We also need a division property.

\begin{lemma}\label{eta:lem:division}
If $F(s)$ is a holomorphic family of order $\mu$ on $\Omega$ and $F(s_0)=0$ for some
$s_0\in\Omega$, then $F(s)=(s-s_0)F_1(s)$ with $F_1$ a holomorphic family of order
$\mu$.
\end{lemma}

\begin{proof} Off the diagonal $\chi K_{F(s)}$ is holomorphic and vanishes at $s_0$,
so is divisible by $s-s_0$. Near the diagonal $\phi F(s)\psi=\Op(\rho^{-\mu(s)}b(s))
+R(s)$, and $\rho^{-\mu(s_0)}b(s_0)$ is a symbol of order $-\infty$, so $b(s_0)$
vanishes to infinite order at the boundary; subtracting
$\rho^{-\mu(s)}\rho^{\mu(s_0)}b(s_0)$ (which changes $R$ by a holomorphic smoothing
family vanishing at $s_0$) we may assume $b(s_0)=0$ and then $b(s)=(s-s_0)b_1(s)$
with $b_1$ smooth and holomorphic. Similarly $R(s)=(s-s_0)R_1(s)$.
\end{proof}

\begin{proposition}\label{eta:prop:admissible}
There is an entire family $E(s)\in\Psido{s}(Z;E)$, $s\in\C$, with
\[
E(0)=\Id
\]
and principal symbol $\abs\xi_g^s\Id_E$ for a Riemannian metric $g$; in particular
$E(s)$ is elliptic for every $s$.
\end{proposition}

\begin{proof} Choose a finite cover by coordinate charts $U_i$ over which $E$ is
trivialized, a subordinate partition of unity $\phi_i$, functions
$\psi_i\in\CIc(U_i)$ equal to $1$ near $\supp\phi_i$, and a metric $g$. Set
\[
E(s)=\sum_i\phi_i\Op\bigl((1+\abs\xi_{g(z)}^2)^{s/2}\Id\bigr)\psi_i.
\]
This is an entire family of order $s$ with principal symbol $\abs\xi_g^s\Id$, and since
$\Op(1)=\Id$, $E(0)=\sum_i\phi_i\psi_i=\sum_i\phi_i=\Id$.
\end{proof}

Using Seeley's complex powers \cite{Seeley} one can also take $E(s)=(\Delta+1)^{s/2}$,
with $\Delta$ a Laplacian, which is in addition invertible for all $s$; we shall not
need this. We call an entire family $E(s)$ of order $s$ with $E(0)=\Id$
\emph{admissible}.

\begin{theorem}[Seeley]\label{eta:thm:seeley}
Let $A(s)$ be a holomorphic family of order $m+s$, $m\in\C$, on a connected open set
$\Omega\subset\C$ which meets the half-plane $\{\Re s<-d-\Re m\}$. Then the function
$\Tr(A(s))$, holomorphic on $\Omega\cap\{\Re s<-d-\Re m\}$, extends to a meromorphic
function on $\Omega$ with at most simple poles, located in $-d-m+\N_0$. The operator $A(s_j)$ at $s_j=j-d-m$ has integral order $j-d$, and the residue there is
\begin{equation}\label{eta:eq:seeley-residue}
\operatorname*{Res}_{s=s_j}\Tr(A(s))=-\res\bigl(A(s_j)\bigr),
\end{equation}
where for $B\in\Psido{k}(Z;E)$, $k\in\Z$,
\[
\res(B)=(2\pi)^{-d}\sum_i\int_{U_i}\phi_i(z)\int_{\abs\xi=1}\tr b^{(i)}_{-d}(z,\xi)\,dS(\xi)\,dz,
\]
$b^{(i)}_{-d}$ being the term of homogeneity $-d$ in the symbol of $B$ in the chart $U_i$
(and $0$ if $k<-d$); here $\{\phi_i\}$ is a partition of unity subordinate to charts
$U_i$ over which $E$ is trivial.
\end{theorem}

\begin{proof} As in the proof of \cref{eta:lem:trace-class},
$A(s)=\sum_i\phi_iA(s)\psi_i+R(s)$ where $R(s)$ is a holomorphic family of
smoothing operators, so $\Tr(R(s))$ is entire. By the definition, $\phi_iA(s)\psi_i
=\Op(\rho^{-m-s}b_i(s))$ modulo holomorphic smoothing families, and we may choose
$\rho=\abs\xi^{-1}$ in $\abs\xi\ge1$. By \eqref{eta:eq:local-trace}, for
$\Re s<-d-\Re m$,
\begin{multline*}
\Tr(\phi_iA(s)\psi_i)=(2\pi)^{-d}\int_U\int_{\abs\xi\le1}\tr(\cdots)\,d\xi\,dz\\
+(2\pi)^{-d}\int_U\int_{\Sphere^{d-1}}\int_0^1 r^{-m-s-d-1}\tr b_i(s,z,r,\omega)\,
dr\,d\omega\,dz,
\end{multline*}
where we have used polar coordinates $\xi=r^{-1}\omega$ in $\abs\xi\ge1$, so that
$d\xi=r^{-d-1}dr\,d\omega$, and $b_i$ is smooth in $r\in[0,1]$ (this is the meaning of
smoothness on $\RC{\R^d}$). The first term is entire. Taylor expanding,
$b_i(s,z,r,\omega)=\sum_{j<J}b_{i,j}(s,z,\omega)r^j+r^Jb_{i,J}(s,z,r,\omega)$ with all
coefficients smooth and holomorphic in $s$, so
\[
\int_0^1r^{-m-s-d-1}b_i\,dr=\sum_{j<J}\frac{b_{i,j}(s,z,\omega)}{j-d-m-s}
+\int_0^1r^{J-m-s-d-1}b_{i,J}\,dr,
\]
and the last integral is holomorphic in $\Re s<J-d-\Re m$. Since $J$ is arbitrary
this gives the meromorphic continuation, with at most simple poles at the points
$s_j$, and residue $-(2\pi)^{-d}\sum_i\int\int\tr b_{i,j}(s_j,z,\omega)\,d\omega\,dz$
there. Now $r^{-m-s_j+j}b_{i,j}(s_j,z,\omega)=\abs\xi^{-d}b_{i,j}(s_j,z,\omega)$ is
precisely the term of homogeneity $-d$ in the symbol of $\phi_iA(s_j)\psi_i$. Finally
the full symbol of $\phi_iB\psi_i$ is $\phi_i$ times the local symbol of $B$ modulo
$S^{-\infty}$, because $\psi_i\equiv1$ near $\supp\phi_i$; this gives
\eqref{eta:eq:seeley-residue}.
\end{proof}

\section{The residue trace and the regularized trace}\label{eta:sec:residue}

\begin{proposition}\label{eta:prop:residue}
For $A\in\Psido{m}(Z;E)$, $m\in\Z$, and any holomorphic family $A(s)$ of order $m+s$
on a connected open set containing $0$ and meeting $\{\Re s<-d-m\}$, with
$A(0)=A$,
\[
\operatorname*{Res}_{s=0}\Tr(A(s))=-\res(A).
\]
The \emph{residue trace} $\TrR(A):=\operatorname*{Res}_{s=0}\Tr(AE(s))$, $E$ admissible,
therefore does not depend on $E$, and
\begin{equation}\label{eta:eq:residue-formula}
\TrR(A)=-\res(A)=-(2\pi)^{-d}\int_{S^*Z}\tr\sigma_{-d}(A),
\end{equation}
where $\sigma_{-d}(A)$ is the homogeneous term of degree $-d$ in the symbol of $A$ and the
integral means, in local coordinates, $\int\int_{\abs\xi=1}\tr\sigma_{-d}(A)\,dS(\xi)\,dz$;
the integrand is a well-defined density on $S^*Z$ although $\sigma_{-d}(A)$ itself depends
on the coordinates. Moreover $\TrR$ vanishes on $\Psido{m}(Z;E)$ for $m<-d$ and on
$\Psido{-\infty}(Z;E)$, and
\[
\TrR([A,B])=0\quad\forall\ A,B\in\Psido{\Z}(Z;E).
\]
\end{proposition}

\begin{proof} The first statement is \eqref{eta:eq:seeley-residue} at $s_j=0$, since the
expression for $\res$ only involves $A(0)=A$. By \cref{eta:lem:hol-composition},
$AE(s)$ is such a family. The formula \eqref{eta:eq:seeley-residue} is a priori written
with a partition of unity; applying it to an operator supported in the intersection of
two charts shows that the density $\bigl(\int_{\abs\xi=1}\tr a_{-d}(z,\xi)dS(\xi)\bigr)dz$
is independent of the coordinates (one can also check this directly, since for a
function $h$ homogeneous of degree $-d$ the form $\iota_R(h\,d\xi)$, $R$ the radial
vector field, is closed on $\R^d\setminus0$). If $m<-d$, $AE(s)$ has order
$m+s$ with $\Re(m+s)<-d$ near $s=0$, so there is no pole.

For the trace property, $\Tr(ABE(s))=\Tr(BE(s)A)$ for $\Re s$ very negative by
\cref{eta:lem:trace-property}, hence for all $s$ by meromorphic continuation, so
\[
\Tr([A,B]E(s))=\Tr\bigl(B[E(s),A]\bigr).
\]
The family $[E(s),A]$ is holomorphic of order $m_A+s$ and vanishes at $s=0$, so by
\cref{eta:lem:division} $[E(s),A]=sC(s)$ with $C$ holomorphic. Since $\Tr(BC(s))$ has
at most a simple pole at $0$, $s\Tr(BC(s))$ is regular there, so $\TrR([A,B])=0$.
\end{proof}

Thus, with our convention that $E(s)$ has order $+s$, the residue trace is minus the
Wodzicki--Guillemin residue \cite{Wodzicki-residue,Guillemin-residue}.

\begin{definition}\label{eta:def:regularized-trace}
For an admissible family $E$ the \emph{regularized trace} of $A\in\Psido{\Z}(Z;E)$ is
the finite part
\[
\Trbar_E(A)=\FP_{s=0}\Tr(AE(s))=\lim_{s\to0}\Bigl(\Tr(AE(s))-\frac1s\TrR(A)\Bigr).
\]
\end{definition}

Clearly $\Trbar_E$ is linear and $\Trbar_E(A)=\Tr(A)$ if $A\in\Psido{m}$, $m<-d$, in
particular on smoothing operators. From the proof of \cref{eta:thm:seeley} the
Laurent coefficients of $\Tr(AE(s))$ at $s=0$ are given by convergent integrals of
the (Taylor coefficients of the) symbol of $AE(s)$, so $\Trbar_E$ is continuous on
each $\Psido{m}(Z;E)$. However $\Trbar_E$ is not a trace, and it depends on $E$.

\begin{lemma}\label{eta:lem:reg-dependence}
If $E$ and $E'$ are admissible then $E'(s)=E(s)+sF(s)$ with $F$ an entire family of
order $s$; setting $D(E,E')=F(0)\in\Psido{0}(Z;E)$,
\[
\Trbar_{E'}(A)=\Trbar_E(A)+\TrR\bigl(AD(E,E')\bigr)\quad\forall\ A\in\Psido{\Z}(Z;E).
\]
\end{lemma}

\begin{proof} The first statement is \cref{eta:lem:division}. Then
$\Tr(AE'(s))=\Tr(AE(s))+s\Tr(AF(s))$; here $AF(s)$ is holomorphic of order $m+s$
with value $AD(E,E')$ at $0$, so $s\Tr(AF(s))$ is regular at $0$ with value
$\TrR(AD(E,E'))$, while the residues of $\Tr(AE'(s))$ and $\Tr(AE(s))$ agree.
\end{proof}

\begin{proposition}\label{eta:prop:affine}
The regularized traces $\Trbar_E$, as $E$ ranges over admissible families, form an
affine space modelled on $\Psido{0}(Z;E)/\Psido{-\infty}(Z;E)$: the map
$E'\longmapsto[D(E,E')]$ identifies $\{\Trbar_{E'}\}$ with
$\Psido{0}/\Psido{-\infty}$.
\end{proposition}

\begin{proof} Every $D\in\Psido{0}$ occurs: $E'(s)=E(s)+sDE(s)$ is admissible with
$D(E,E')=D$. By \cref{eta:lem:reg-dependence} it remains to see that the functional
$A\mapsto\TrR(AD)$ on $\Psido{\Z}$ vanishes if and only if $D\in\Psido{-\infty}$. If
$D$ is smoothing so is $AD$. Conversely, if $D\notin\Psido{-\infty}$ let $k\in\N_0$
be such that $D\in\Psido{-k}$ has non-vanishing principal symbol $\sigma_{-k}(D)$. For
$A\in\Psido{k-d}$ the term of homogeneity $-d$ in the symbol of $AD$ is
$\sigma_{k-d}(A)\sigma_{-k}(D)$, so by \eqref{eta:eq:residue-formula}
\[
\TrR(AD)=-(2\pi)^{-d}\int_{S^*Z}\tr\bigl(\sigma_{k-d}(A)\sigma_{-k}(D)\bigr).
\]
Choosing $\sigma_{k-d}(A)=\chi\,\sigma_{-k}(D)^*$ (extended homogeneously) with
$0\le\chi$, $\chi>0$ somewhere where $\sigma_{-k}(D)\ne0$, makes this non-zero.
\end{proof}

\subsection{The trace-defect formula} The failure of $\Trbar_E$ to be a trace is
measured by a derivation.

\begin{definition}\label{eta:def:derivation}
For an admissible family $E$ and $A\in\Psido{\Z}(Z;E)$ let
\[
D_EA=\frac{d}{ds}[E(s),A]\Big|_{s=0},
\]
so that $[E(s),A]=sC(s)$ with $C$ holomorphic of order $m+s$ and $C(0)=D_EA$.
\end{definition}

Formally, if $E(s)=Q^s$ then $D_EA=[\log Q,A]$; $\log Q$ is not in the algebra, but
the commutator is.

\begin{proposition}\label{eta:prop:derivation}
$D_E$ is a derivation of $\Psido{\Z}(Z;E)$ mapping $\Psido{m}$ into $\Psido{m}$; if the
principal symbol of $E(s)$ is $\abs\xi_g^s\Id$ (as in \cref{eta:prop:admissible}) then
$D_E(\Psido{m})\subset\Psido{m-1}$ and
\begin{equation}\label{eta:eq:DE-symbol}
\sigma_{m-1}(D_EA)=-i\{\log\abs\xi_g,\sigma_m(A)\}.
\end{equation}
For two admissible families, $D_{E'}A=D_EA+[D(E,E'),A]$. Moreover
\[
\TrR(D_EA)=0\quad\forall\ A\in\Psido{\Z}(Z;E).
\]
\end{proposition}

\begin{proof} Differentiating $[E(s),AB]=[E(s),A]B+A[E(s),B]$ at $s=0$ gives the
derivation property, and the difference formula by
differentiating $[E(s)+sF(s),A]$. If the principal symbol $e_s$ of $E(s)$ is
scalar, the principal symbol of $[E(s),A]$ of order $m+s$ vanishes, and the term of
order $m+s-1$ is
\[
-i\{e_s,\sigma_m(A)\}=-is\abs\xi_g^s\{\log\abs\xi_g,\sigma_m(A)\};
\]
dividing by $s$ and setting $s=0$ gives \eqref{eta:eq:DE-symbol}. Finally, the
family $C(s)$ has order $m+s$ and value $D_EA$ at $0$, so $\TrR(D_EA)$ is the residue
of $\Tr(C(s))=s^{-1}\Tr([E(s),A])$; but $\Tr([E(s),A])=0$ for $\Re s$ very negative by
\cref{eta:lem:trace-property}, hence identically.
\end{proof}

One can show that $D_E$ is not an inner derivation (see \cref{eta:ex:not-inner} for the
circle).

\begin{lemma}[Trace-defect formula]\label{eta:lem:trace-defect}
For all $A,B\in\Psido{\Z}(Z;E)$,
\[
\Trbar_E([A,B])=\TrR\bigl((D_EA)B\bigr)=\TrR\bigl(BD_EA\bigr).
\]
\end{lemma}

\begin{proof} As in the proof of \cref{eta:prop:residue},
$\Tr([A,B]E(s))=\Tr(B[E(s),A])=s\Tr(BC(s))$, which is regular at $s=0$ with value
$\operatorname*{Res}_{s=0}\Tr(BC(s))=\TrR(BD_EA)$. Since $\TrR([A,B])=0$, this value is
$\Trbar_E([A,B])$; the second equality is the trace property of $\TrR$.
\end{proof}

We shall apply this to operator-valued differential forms. If $X$ is a manifold and
$\alpha=\sum_IA_I\,dx^I$, $\beta=\sum_JB_J\,dx^J$ are smooth forms of degrees $p$ and
$q$ with values in $\Psido{\Z}(Z;E)$, products are taken in the algebra and in the
exterior algebra, $\Trbar_E$, $\TrR$ and $D_E$ act on coefficients, and the graded
commutator is $[\alpha,\beta]=\alpha\beta-(-1)^{pq}\beta\alpha
=\sum_{I,J}[A_I,B_J]\,dx^I\wedge dx^J$. Hence:

\begin{corollary}\label{eta:cor:graded-defect}
For operator-valued forms $\alpha$, $\beta$ as above,
$\Trbar_E([\alpha,\beta])=\TrR\bigl((D_E\alpha)\wedge\beta\bigr)$, and $\TrR$ vanishes on
graded commutators. Moreover $d\,\Trbar_E(\alpha)=\Trbar_E(d\alpha)$, by continuity
and linearity of $\Trbar_E$.
\end{corollary}

\section{The circle}\label{eta:sec:circle}

On $\Sphere=\R/2\pi\Z$, with variable $\theta$ and dual variable $\xi$, we use the
admissible family
\begin{equation}\label{eta:eq:circle-E}
E(s)e^{ik\theta}=\lambda_k^se^{ik\theta},\qquad\lambda_k=(1+k^2)^{1/2},
\end{equation}
acting diagonally on $\CI(\Sphere;\C^N)$. It is the Fourier multiplier with symbol
$\ang\xi^s$, an entire family of invertible elliptic operators of order $s$ with
$E(0)=\Id$. From now on $\Trbar=\Trbar_E$ and $D_E$ refer to this family. Here
$S^*\Sphere=\Sphere\times\{\pm1\}$.

\begin{lemma}\label{eta:lem:circle}
\begin{enumerate}[label=(\roman*)]
\item For $A\in\Psido{-1}(\Sphere;\C^N)$ with symbol term $a_{-1}$ of homogeneity $-1$,
\[
\TrR(A)=-\frac1{2\pi}\int_0^{2\pi}\tr\bigl(a_{-1}(\theta,1)+a_{-1}(\theta,-1)\bigr)\,d\theta.
\]
\item For $A\in\Psido{m}(\Sphere;\C^N)$ with principal symbol $a$, $D_EA\in\Psido{m-1}$
has principal symbol
\[
\sigma_{m-1}(D_EA)(\theta,\xi)=-i\,\xi^{-1}\pa_\theta a(\theta,\xi);
\]
in fact $D_EA$ has full symbol $\sum_{k\ge1}\frac1{k!}\pa_\xi^k(\log\ang\xi)D_\theta^ka$.
\item If $f\in\CI(\Sphere;M(N,\C))$ acts by multiplication and $\Pi$ is the Szeg\H{o}
projector then
\[
\Trbar(f)=0,\qquad\Trbar(\Pi f\Pi)=\frac1{4\pi}\int_0^{2\pi}\tr f\,d\theta;
\]
in particular $\Trbar(\Pi)=\frac12$, and $\TrR(f)=\TrR(\Pi f\Pi)=0$.
\end{enumerate}
\end{lemma}

\begin{proof} (i) is \eqref{eta:eq:residue-formula} for $d=1$. For (ii), since the
symbol $e_s=\ang\xi^s$ depends only on $\xi$, $\Op(a)\Op(e_s)=\Op(ae_s)$ while
$\Op(e_s)\Op(a)$ has symbol $\sum_k\frac1{k!}\pa_\xi^ke_sD_\theta^ka$; every term with
$k\ge1$ is divisible by $s$, with $\frac{d}{ds}\pa_\xi^ke_s|_{s=0}=\pa_\xi^k\log\ang\xi$.
The leading term is $\pa_\xi\log\ang\xi\,D_\theta a=\xi\ang\xi^{-2}(-i\pa_\theta a)$.
For (iii), the multiplication operator has diagonal matrix entries
$\langle fe_k,e_k\rangle=\frac1{2\pi}\int f\,d\theta$ for every $k\in\Z$ (in
$\Pi f\Pi$ only for $k\ge0$), so for $\Re s<-1$
\[
\Tr(fE(s))=\frac1{2\pi}\int\tr f\,d\theta\sum_{k\in\Z}\lambda_k^s,\quad
\Tr(\Pi f\Pi E(s))=\frac1{2\pi}\int\tr f\,d\theta\sum_{k\ge0}\lambda_k^s.
\]
Now $\sum_{k\ge0}\lambda_k^s=1+\zeta(-s)+h(s)$ with
$h(s)=\sum_{k\ge1}(\lambda_k^s-k^s)$; since $\lambda_k^s-k^s=O(\abs s k^{\Re s-2})$,
$h$ is holomorphic in $\Re s<1$ with $h(0)=0$. The Riemann zeta function is regular
at $0$ with $\zeta(0)=-\frac12$, so the finite part (in fact the value) at $s=0$ is
$\frac12$, and $\sum_{k\in\Z}\lambda_k^s=2\sum_{k\ge0}\lambda_k^s-1$ has value $0$.
\end{proof}

\begin{example}\label{eta:ex:defect-check}
Let $A$ be multiplication by $e^{i\theta}$ and $B=e^{-i\theta}\Pi$. Then $AB=\Pi$ and
$BA=e^{-i\theta}\Pi e^{i\theta}$ is the projection onto the span of the $e^{ik\theta}$,
$k\ge-1$, so $[A,B]=-\pi_{-1}$, minus the projection onto $e^{-i\theta}$, and $\Trbar([A,B])=-1$. On the other hand, by \cref{eta:lem:circle}(ii),
$(D_EA)B$ has symbol $\xi^{-1}$ in $\xi>0$ and is rapidly decreasing in $\xi<0$, so
$\TrR((D_EA)B)=-\frac1{2\pi}\int_0^{2\pi}1\,d\theta=-1$, as it must be by
\cref{eta:lem:trace-defect}. This confirms the signs in the trace-defect formula.
\end{example}

\section{Odd Chern forms and the Toeplitz index}\label{eta:sec:chern}

\subsection{Toeplitz operators} We recall from \cref{ch:index}. The Szeg\H{o} projector
$\Pi\in\Psido{0}(\Sphere)$ projects onto the span $\Hardy$ of the $e^{ik\theta}$,
$k\ge0$; its symbol is $1$ on $\xi>0$ and its full symbol is rapidly decreasing as
$\xi\to-\infty$. The Toeplitz algebra
\[
\PsiT{m}(\Sphere;\C^N)=\Pi\Psido{m}(\Sphere;\C^N)\Pi
\]
acts on $\Hardy\otimes\C^N$, with identity $\Pi$; the full symbols of its elements are
rapidly decreasing as $\xi\to-\infty$, and we write $\sigma_0(A)(\theta)$ for the
principal symbol at $(\theta,+1)$. An element of $\PsiT{0}$ is Fredholm on
$\Hardy\otimes\C^N$ if and only if $\sigma_0(A)$ is invertible, and then
\begin{equation}\label{eta:eq:toeplitz-index}
\ind(A)=\frac{i}{2\pi}\int_0^{2\pi}\tr\bigl(\sigma_0(A)^{-1}\pa_\theta\sigma_0(A)\bigr)d\theta
=-\wind\det\sigma_0(A).
\end{equation}
For instance the shift $U=\Pi e^{i\theta}\Pi$ has index $-1$. Let
\[
\GT(\Sphere;\C^N)=\{B\in\PsiT{0}(\Sphere;\C^N);\ B\text{ invertible on }\Hardy\otimes\C^N,
\ B^{-1}\in\PsiT{0}(\Sphere;\C^N)\}
\]
and $G^{-\infty}_{\mathcal T}=\GT\cap(\Pi+\PsiT{-\infty})$. Symbols of elements of
$\GT$ are loops in $\GL(N,\C)$ with winding number $0$. Products of operators on
$\Hardy\otimes\C^N$ are products in $\Psido{\Z}(\Sphere;\C^N)$, so $\Trbar$ and $\TrR$
apply to elements of $\PsiT{\Z}$.

\subsection{Odd Chern forms} For a manifold $W$ and $a\in\CI(W;\GL(N,\C))$ set
\begin{equation}\label{eta:eq:chodd}
\Chodd(a)=\sum_{k\ge0}c_k\tr\bigl((a^{-1}da)^{2k+1}\bigr),\qquad
c_k=-\frac{k!}{(2\pi i)^{k+1}(2k+1)!}.
\end{equation}
The same formula, with $\tr$ replaced by $\Tr$, defines $\Chodd$ for maps into the
group $G^{-\infty}$ of \cref{ch:index}. The forms are closed: with
$\nu=a^{-1}da$ we have $d\nu=-\nu^2$, hence
\begin{equation}\label{eta:eq:d-power}
d(\nu^{2k+1})=\sum_{j=0}^{2k}(-1)^j\nu^jd\nu\,\nu^{2k-j}=-\nu^{2k+2},
\end{equation}
and $\tr(\nu^{2k+2})=0$ because $\tr(\nu\wedge\nu^{2k+1})=-\tr(\nu^{2k+1}\wedge\nu)$. The
constants $c_k$ are fixed by \cref{eta:prop:toeplitz-chern} below; note
$c_0=i/2\pi$, consistent with \eqref{eta:eq:toeplitz-index}.

\emph{Conventions.} On a product $W\times\Sphere$ every form can be written uniquely as
$\omega=\omega_0+\omega_1\wedge d\theta$ with $\omega_0,\omega_1$ not containing
$d\theta$, and we set $\int_\Sphere\omega=\int_0^{2\pi}\omega_1\,d\theta$, a form on
$W$. With this (`fibre last') convention $d\int_\Sphere=\int_\Sphere d$,
$\int_\Sphere(\pi^*\alpha\wedge\omega)=\alpha\wedge\int_\Sphere\omega$, and
$\int_{W\times\Sphere}=\int_W\int_\Sphere$ when $W\times\Sphere$ carries the product
orientation (first $W$, then $d\theta$).

\begin{lemma}\label{eta:lem:transgression}
Let $a_t\in\CI(W;\GL(N,\C))$ depend smoothly on $t\in[0,1]$ and regard $a$ as a map on
$W\times[0,1]$. Then
\[
\Chodd(a)=\Chodd(a_t)+\Theta(a_t)\wedge dt,\qquad
\Theta(a_t)=\sum_{k\ge0}(2k+1)c_k\tr\bigl((a_t^{-1}da_t)^{2k}a_t^{-1}\pa_ta_t\bigr),
\]
and $\pa_t\Chodd(a_t)=d\Theta(a_t)$.
\end{lemma}

\begin{proof} Write $a^{-1}Da=\nu+\lambda\,dt$, $D=d+dt\,\pa_t$, $\nu=a_t^{-1}da_t$,
$\lambda=a_t^{-1}\pa_ta_t$. The $dt$-component of $(\nu+\lambda dt)^{2k+1}$ is
\[
\sum_j\nu^j\lambda\,dt\,\nu^{2k-j}=\sum_j(-1)^{2k-j}\nu^j\lambda\nu^{2k-j}\wedge dt,
\]
and $\tr(\nu^j\lambda\nu^{2k-j})=(-1)^{j(2k-j)}\tr(\lambda\nu^{2k})=(-1)^j\tr(\lambda\nu^{2k})$,
giving $(2k+1)\tr(\nu^{2k}\lambda)\wedge dt$. The second identity is the
$dt$-component of $D\Chodd(a)=0$, using $dt\wedge\pa_t\omega=-\pa_t\omega\wedge dt$
for $\omega$ of odd degree.
\end{proof}

Applying this to the rotation homotopy from $g\oplus h$ to $gh\oplus\Id$ shows that
$\Chodd(gh)-\Chodd(g)-\Chodd(h)$ is exact (\cref{eta:ex:chodd-product}).

\begin{proposition}\label{eta:prop:toeplitz-chern}
Let $W$ be compact, $b\in\CI(W\times\Sphere;\GL(N,\C))$ and let $T(w)\in\PsiT{0}(\Sphere;\C^N)$
be a smooth family with symbol $b(w,\cdot)$. Then the Chern character of the index
bundle $\ind(T)\in K^0(W)$ is
\begin{equation}\label{eta:eq:toeplitz-chern}
\Ch(\ind T)=\Bigl[\int_\Sphere\Chodd(b)\Bigr]\in H^{\mathrm{even}}(W;\C),
\end{equation}
where $\Ch$ of a bundle with curvature $\Omega$ is represented by $\tr\exp(i\Omega/2\pi)$.
\end{proposition}

\begin{proof} This is the families index theorem for Toeplitz operators of
\cref{ch:families}; we recall the structure of the argument and verify the
constants. Both sides are additive under direct sums, unchanged by stabilization and
homotopy invariant (the right side by \cref{eta:lem:transgression}). Multiplying
$T$ by the $\theta$-independent invertible matrix $b(w,1)^{-1}$ changes neither side
(for the right side use \cref{eta:ex:chodd-product} and the fact that
$\int_\Sphere\Chodd(g)=0$ for $\theta$-independent $g$). By the homotopy results of
\cref{ch:bott}, after these reductions and stabilization $b$ may be deformed to
$\pi_-e^{-i\theta}+\pi_0+\pi_+e^{i\theta}$, with $\pi_\pm,\pi_0$ smooth families of
commuting projections on some $\C^M$ summing to $\Id$. This is the product of the
commuting symbols $\pi_+e^{i\theta}+(\Id-\pi_+)$ and $\pi_-e^{-i\theta}+(\Id-\pi_-)$;
the index is additive under composition (the Toeplitz operator of a product differs
from the product of the Toeplitz operators by a compact family) and $\Chodd$ is additive
up to exact forms (\cref{eta:ex:chodd-product}), so it suffices to consider the two
factors. The second is obtained from a symbol of the first type by $\theta\mapsto-\theta$,
which reverses the sign of both sides. So it suffices to consider $b=\pi e^{i\theta}+\pi'$, with
Toeplitz operator $\pi U+\pi'$: it is injective with cokernel $\pi\C^M$ (in the
constants), so $\ind(T)=-[V]$, $V=\Ran\pi$. The connection $\pi d$ on $V$ has
curvature $\Omega=\pi\,d\pi\wedge d\pi\,\pi$, so the left side of
\eqref{eta:eq:toeplitz-chern} is $-\sum_k\frac1{k!}(\frac i{2\pi})^k\tr\Omega^k$.

On the other hand, writing $d'$ for the differential on $W\times\Sphere$,
\[
b^{-1}d'b=i\pi\,d\theta+\nu,\quad \nu=f\,\pi\,d\pi\,\pi'+g\,\pi'd\pi\,\pi,\quad
f=1-e^{-i\theta},\ g=e^{i\theta}-1,
\]
using $\pi\,d\pi\,\pi=0$ and $\pi'd\pi\,\pi'=0$. The $d\theta$-component of
$(b^{-1}d'b)^{2k+1}$ has trace $(2k+1)i\tr(\pi\nu^{2k}\pi)\wedge d\theta$ (by
cyclicity, as in \cref{eta:lem:transgression}) and $\pi\nu^{2k}\pi=(fg)^k
(\pi\,d\pi\,\pi'd\pi\,\pi)^k=(-1)^k(2-2\cos\theta)^k\Omega^k$. Since
$\int_0^{2\pi}(2-2\cos\theta)^kd\theta=2\pi\frac{(2k)!}{(k!)^2}$,
\[
\int_\Sphere\tr\bigl((b^{-1}d'b)^{2k+1}\bigr)=2\pi i(-1)^k\frac{(2k+1)!}{(k!)^2}\tr\Omega^k,
\]
and multiplying by $c_k$ gives exactly $-\frac1{k!}(\frac i{2\pi})^k\tr\Omega^k$.
\end{proof}

\subsection{The Fredholm determinant} Recall from \cref{ch:index} that on $G^{-\infty}$
the Fredholm determinant satisfies $d\log\det A=\Tr(A^{-1}dA)$, so that the degree-one
part $c_0\Tr(A^{-1}dA)$ of $\Chodd$ is $\frac{i}{2\pi}d\log\det A$, and that
$\frac1{2\pi i}\int_\gamma\Tr(A^{-1}dA)\in\Z$ for every loop $\gamma$; this integer is the
case $W=\{\mathrm{pt}\}$ of \eqref{eta:eq:toeplitz-index}--\eqref{eta:eq:toeplitz-chern}.

\section{Eta forms}\label{eta:sec:eta}

Let $X$ be a manifold and $B\in\CI(X;\GT(\Sphere;\C^N))$ a smooth family of invertible
Toeplitz operators. Then
\[
\alpha=B^{-1}dB\in\CI(X;\Lambda^1\otimes\PsiT{0}(\Sphere;\C^N))
\]
and the symbol $b=\sigma_0(B)\in\CI(X\times\Sphere;\GL(N,\C))$ satisfies $\wind\det b(x,\cdot)=0$.
We write
\[
\nu=b^{-1}d_Xb,\qquad\mu=b^{-1}\pa_\theta b,
\]
for the components of $b^{-1}d'b=\nu+\mu\,d\theta$ on $X\times\Sphere$.

\begin{definition}\label{eta:def:eta}
The \emph{eta forms} of $B$ are
\[
\eta_{2k+1}(B)=\Trbar\bigl((B^{-1}dB)^{2k+1}\bigr),\qquad\eta(B)=\sum_{k\ge0}c_k\eta_{2k+1}(B),
\]
and the associated \emph{symbolic correction forms} are
\[
\gamma_{2k+1}(b)=\frac{i}{4\pi}\int_0^{2\pi}\tr\bigl(\nu^{2k+1}\mu\bigr)d\theta,\qquad
\gamma(b)=\sum_{k\ge0}c_k\gamma_{2k+1}(b).
\]
\end{definition}

The $\gamma$-forms are local in the symbol $b$; they vanish if $b$ is independent of
$\theta$. If $B\in\CI(X;G^{-\infty}_{\mathcal T})$ then $B^{-1}dB$ is smoothing, so
$\Trbar=\Tr$ and $\eta(B)=\Chodd(B)$, while $b=\Id$ and $\gamma(b)=0$. Thus $\eta$
\emph{extends the odd Chern character from the normal subgroup $G^{-\infty}_{\mathcal T}$
to the whole group} $\GT$, at the expense of no longer being closed.

\begin{theorem}\label{eta:thm:variation}
For every smooth family $B\in\CI(X;\GT(\Sphere;\C^N))$ with symbol $b$,
\begin{equation}\label{eta:eq:variation}
d\bigl(\eta_{2k+1}(B)+\gamma_{2k+1}(b)\bigr)=\frac{i}{4\pi(2k+3)}\int_\Sphere
\tr\bigl((b^{-1}d'b)^{2k+3}\bigr)
\end{equation}
and consequently
\begin{equation}\label{eta:eq:variation-total}
d\bigl(\eta(B)+\gamma(b)\bigr)=-\int_\Sphere\Chodd(b).
\end{equation}
\end{theorem}

\begin{proof} By \cref{eta:cor:graded-defect} and \eqref{eta:eq:d-power} applied to
$\alpha$,
\[
d\eta_{2k+1}(B)=\Trbar\bigl(d(\alpha^{2k+1})\bigr)=-\Trbar(\alpha^{2k+2})
=-\frac12\Trbar\bigl([\alpha,\alpha^{2k+1}]\bigr)
=-\frac12\TrR\bigl((D_E\alpha)\wedge\alpha^{2k+1}\bigr),
\]
since the graded commutator is $[\alpha,\alpha^{2k+1}]=2\alpha^{2k+2}$. The operator
$(D_E\alpha)\wedge\alpha^{2k+1}$ has order $-1$, and its full symbol is rapidly
decreasing as $\xi\to-\infty$ because $\alpha$ is Toeplitz. By
\cref{eta:lem:circle}(i),(ii) only the principal symbol at $\xi=+1$, namely
$-i(\pa_\theta\nu)\nu^{2k+1}$, contributes:
\[
d\eta_{2k+1}(B)=-\frac12\cdot\Bigl(-\frac1{2\pi}\Bigr)\int_0^{2\pi}\tr\bigl(-i\,\pa_\theta\nu\wedge
\nu^{2k+1}\bigr)d\theta=-\frac{i}{4\pi}\int_0^{2\pi}\tr\bigl(\pa_\theta\nu\wedge\nu^{2k+1}\bigr)d\theta.
\]
Differentiating $\nu=b^{-1}d_Xb$ and $\mu=b^{-1}\pa_\theta b$ gives the Maurer--Cartan
identity $\pa_\theta\nu=d_X\mu+\nu\mu-\mu\nu$. By cyclicity
$\tr(\nu\mu\nu^{2k+1})=-\tr(\mu\nu^{2k+2})$, and by \eqref{eta:eq:d-power}
$\tr(d\mu\wedge\nu^{2k+1})=d\tr(\mu\nu^{2k+1})+\tr(\mu\nu^{2k+2})$. Hence
\[
\tr(\pa_\theta\nu\wedge\nu^{2k+1})=d\tr(\mu\nu^{2k+1})-\tr(\mu\nu^{2k+2}),
\]
and, $\mu$ being of degree $0$,
\[
d\eta_{2k+1}(B)=-d\gamma_{2k+1}(b)+\frac{i}{4\pi}\int_0^{2\pi}\tr(\nu^{2k+2}\mu)\,d\theta.
\]
Finally, the $d\theta$-component of $(\nu+\mu\,d\theta)^{2k+3}$ is
$\sum_{j=0}^{2k+2}(-1)^j\nu^j\mu\nu^{2k+2-j}\wedge d\theta$, whose trace is
$(2k+3)\tr(\nu^{2k+2}\mu)\wedge d\theta$ since $\tr(\nu^j\mu\nu^{2k+2-j})=(-1)^j\tr(\mu\nu^{2k+2})$.
This proves \eqref{eta:eq:variation}. For \eqref{eta:eq:variation-total} note that the
constants satisfy
\[
c_k\cdot\frac{i}{4\pi(2k+3)}=-c_{k+1},
\]
so the sum over $k$ of \eqref{eta:eq:variation} gives $-\int_\Sphere\Chodd(b)$ except
for the term $c_0\int_\Sphere\tr(b^{-1}d'b)=c_0\int_0^{2\pi}\tr\mu\,d\theta$ of degree
$0$. By \eqref{eta:eq:toeplitz-index} this is $\ind(B(x))=0$.
\end{proof}

Thus $-\bigl(\eta(B)+\gamma(b)\bigr)$ is an explicit primitive of the form
$\int_\Sphere\Chodd(b)$, which by \cref{eta:prop:toeplitz-chern} represents the Chern
character of the index bundle of any Toeplitz family with symbol $b$, a class that
vanishes here because the family $B$ is invertible. Note that the variation of the eta
form depends only on the leading part of the symbol, not on the lower order terms or
on the smoothing part of $B$. The situation can be summarized in the diagram
\[
\begin{tikzcd}[column sep=small]
\Chodd & \eta\arrow[l,"\text{restrict}"']\arrow[r,"d(\cdot+\gamma)"] & -\int_\Sphere\Chodd\\
G^{-\infty}_{\mathcal T}(\Sphere;\C^N)\arrow[r] & \GT(\Sphere;\C^N)\arrow[r,"\sigma_0"] &
\CI(\Sphere;\GL(N,\C))_0
\end{tikzcd}
\]
where the right-hand group consists of loops of winding number $0$. In the stabilized
setting of \cref{ch:bott} (where $\C^N$ is replaced by $\Psido{-\infty}(Y;F)$ and $\tr$
by the trace on $Y$; all the arguments above go through verbatim) the bottom row is
the classifying sequence for $K$-theory, with weakly contractible middle group, and the
eta form is a transgression form for the even Chern character of the quotient. That
the transgression only involves the principal symbol is a higher analogue of the
`miracle of the loop group' of Pressley and Segal \cite{Pressley-Segal}.

\begin{proposition}\label{eta:prop:eta-det}
Let $B\in\CI(X;\GT(\Sphere;\C^N))$ and $g\in\CI(X;G^{-\infty}_{\mathcal T})$. Then
\begin{enumerate}[label=(\roman*)]
\item $\eta_1(g)=\Tr(g^{-1}dg)=d\log\det g$;
\item $\eta_1(Bg)=\eta_1(B)+d\log\det g$ and $\eta_1(gB)=\eta_1(B)+d\log\det g$;
\item $d\bigl(\eta_1(B)+\gamma_1(b)\bigr)=\dfrac{i}{12\pi}\displaystyle\int_\Sphere\tr\bigl((b^{-1}d'b)^3\bigr)$.
\end{enumerate}
\end{proposition}

\begin{proof} (i) was noted above and (iii) is \eqref{eta:eq:variation} for $k=0$. For
(ii), $(Bg)^{-1}d(Bg)=g^{-1}\alpha g+g^{-1}dg$, and
$g^{-1}\alpha g-\alpha=[g^{-1}\alpha,g]$, so by \cref{eta:lem:trace-defect}
\[
\Trbar(g^{-1}\alpha g)-\Trbar(\alpha)=\TrR\bigl(D_E(g^{-1}\alpha)\,g\bigr).
\]
Since $g-\Pi$ is smoothing and $\sigma_0(g^{-1}\alpha)=\nu$, by
\cref{eta:lem:circle} the right side is
\[
\frac{i}{2\pi}\int_0^{2\pi}\pa_\theta\tr\nu\,d\theta=0.
\] Also $\Trbar(g^{-1}dg)=\Tr(g^{-1}dg)$. For $gB$, we have
$(gB)^{-1}d(gB)=B^{-1}(g^{-1}dg)B+\alpha$ and
$\Tr(B^{-1}(g^{-1}dg)B)=\Tr(g^{-1}dg)$ because $g^{-1}dg$ is of trace class.
\end{proof}

\begin{remark}\label{eta:rem:determinant-connection}
Consider the principal bundle $\GT\longrightarrow\GT/G^{-\infty}_{\mathcal T}$ and the line
bundle $L$ associated to it by the character $\det:G^{-\infty}_{\mathcal T}\longrightarrow\C^*$.
By (ii), and since $\gamma_1(\sigma_0(Bg))=\gamma_1(\sigma_0(B))$, the $1$-form
$\eta_1+\gamma_1$ on $\GT$ is a connection form for $L$ (its restriction to the
fibres is $d\log\det$ and it transforms correctly); by (iii) its curvature is the
$2$-form $\frac{i}{12\pi}\int_\Sphere\tr((b^{-1}d'b)^3)$, which depends only on the
symbol and so descends. The first Chern form $\frac{i}{2\pi}d(\eta_1+\gamma_1)$ of $L$
is therefore $-\frac1{24\pi^2}\int_\Sphere\tr((b^{-1}d'b)^3)$, which is minus the
degree-two part of $\int_\Sphere\Chodd(b)$. This is consistent with the identification
of $L$ with the determinant line bundle $\Lambda^{\max}(\Nul)^*\otimes\Lambda^{\max}
(\text{coker})$ of the index bundle, whose first Chern class is $-c_1(\ind)$. Thus
$\eta_1$ is a regularized form of $d\log\det$ on the Toeplitz group; see also
\cref{eta:sec:determinant}.
\end{remark}

\section{Product-type Toeplitz operators}\label{eta:sec:product-toeplitz}

\subsection{The product-type calculus on $M\times\Sphere$} Let $M$ be a compact
manifold without boundary of dimension $n\ge1$ and consider the trivial fibration
$\pi:M\times\Sphere\longrightarrow M$. We use the following facts from
\cref{ch:product} about the product-type operators $\Psipt{0,0}(M\times\Sphere;\C^N)$
associated to this fibration.

\begin{enumerate}[label=(P\arabic*)]
\item\label{eta:P1} They form an algebra, acting on $\CI(M\times\Sphere;\C^N)$ and
boundedly on $L^2$, closed under adjoints. There are two multiplicative symbol maps.
The \emph{principal symbol} $\sigma(A)$ is a smooth function on the blow-up of
$S^*(M\times\Sphere)$ along $S^*M\times\Sphere$ (the covectors annihilating the fibres,
$\tau=0$ where $\tau$ is dual to $\theta$); this blow-up is the disjoint union of the
two closed half-sphere bundles
\[
S^*_\pm=\{\pm\tau\ge0\}\subset S^*(M\times\Sphere),\qquad\pa S^*_\pm=S^*M\times\Sphere.
\]
The \emph{base family} $\beta(A)\in\CI(S^*M;\Psido{0}(\Sphere;\C^N))$ is a family of
operators on the fibre parametrized by $S^*M$.
\item\label{eta:P2} The symbols are compatible: the principal symbol of $\beta(A)(x)$ on
$S^*\Sphere=\Sphere\times\{\pm1\}$ is the restriction of $\sigma(A)$ to
$\{x\}\times\Sphere\subset\pa S^*_\pm$. Conversely any compatible pair, depending
smoothly on parameters, is the joint symbol of a smooth family of operators.
\item\label{eta:P3} If $\sigma(A)=0$ and $\beta(A)=0$ then $A$ is compact on $L^2$.
\item\label{eta:P4} If $A$ is \emph{fully elliptic}, i.e.\ $\sigma(A)$ and all $\beta(A)(x)$ are
invertible, then $A$ is Fredholm on $L^2$ and on $\CI$, with the same index; the
index depends only on the joint symbol and is constant along smooth families of
fully elliptic operators.
\item\label{eta:P5} The ordinary operators $\Psido{0}(M\times\Sphere;\C^N)$ are product-type,
with $\sigma(P)$ the lift of the usual principal symbol and $\beta(P)(x)$ the
multiplication operator by $\sigma(P)|_{\{x\}\times\Sphere}$. Fibrewise families
$Q\in\CI(M;\Psido{0}(\Sphere;\C^N))$ are product-type, with $\sigma(Q)$ the lift of
the fibre symbol and $\beta(Q)(m,\xi)=Q(m)$; in particular multiplication by
$h\in\CI(M\times\Sphere;M(N,\C))$ has symbol $h$ and base family the multiplication by
$h(m,\cdot)$.
\end{enumerate}

Since the Szeg\H{o} projector $\Pi$, acting in $\theta$, is a fibrewise operator, it is in
$\Psipt{0,0}$ with $\sigma(\Pi)=1$ on $S^*_+$, $\sigma(\Pi)=0$ on $S^*_-$, and
$\beta(\Pi)=\Pi$. The \emph{product-type Toeplitz algebra} is
\[
\PsiptT{0,0}(M\times\Sphere;\C^N)=\Pi\,\Psipt{0,0}(M\times\Sphere;\C^N)\,\Pi,
\]
acting on $\Hardy(M\times\Sphere;\C^N)=\Pi\CI(M\times\Sphere;\C^N)$ (or its $L^2$
closure), with identity $\Pi$. For $A$ in this algebra the principal symbol vanishes
on $S^*_-$; we set
\begin{gather*}
a=\sigma(A)|_{S^*_+}\in\CI(S^*_+;M(N,\C)),\\
\beta=\beta(A)\in\CI(S^*M;\PsiT{0}(\Sphere;\C^N)),\quad b=a|_{\pa S^*_+}=\sigma_0(\beta).
\end{gather*}
$A$ is \emph{fully elliptic} if $a$ takes values in $\GL(N,\C)$ and
$\beta\in\CI(S^*M;\GT(\Sphere;\C^N))$; then $A$ is Fredholm on $\Hardy$ (apply
\ref{eta:P4} to $A+\Id-\Pi$) and its index has the properties in \ref{eta:P4}. We
write $\ind(A)$ for the index on $\Hardy$.

\subsection{Geometry and orientations} Fix a metric on $M$ and use the product metric
on $M\times\Sphere$. In each fibre of $S^*_+\to M\times\Sphere$, stereographic
projection $(\xi,\tau)\mapsto\xi/(1+\tau)$ identifies the closed half-sphere
$\{\abs\xi^2+\tau^2=1,\ \tau\ge0\}$ diffeomorphically with the closed unit ball in
$T^*_mM$; thus $S^*_+\cong\overline{B^*M}\times\Sphere$ is a ball bundle over
$M\times\Sphere$, with `centre' the section $(\xi,\tau)=(0,1)$, i.e.\ $d\theta$.

We orient cotangent bundles by the Liouville form: if $\omega=\sum_jd\xi_j\wedge dx_j$
then $T^*M$ is oriented by $\omega^n/n!$. The ball bundle $\overline{B^*M}$ inherits
this orientation, and $S^*M=\pa\overline{B^*M}$ is given the boundary orientation
(outward normal first); we call this the \emph{Liouville orientation} of $S^*M$, and
similarly for $S^*(M\times\Sphere)$, whose open subset $\{\tau>0\}$ orients $S^*_+$.
For $n=1$ the two components $\{\xi=\pm1\}$ of $S^*M\cong M\times\{\pm1\}$ receive
opposite orientations ($\pm dx$ if $M$ is oriented by $dx$).

\begin{lemma}\label{eta:lem:orientation}
With these orientations, $\pa S^*_+=S^*M\times\Sphere$ with the product orientation,
$S^*M$ first and $d\theta$ last.
\end{lemma}

\begin{proof} The Liouville form of $T^*(M\times\Sphere)$ is $\omega_M+d\tau\wedge d\theta$,
and $(\omega_M+d\tau\wedge d\theta)^{n+1}/(n+1)!=\omega_M^n/n!\wedge d\tau\wedge d\theta$.
At a point with $\tau=0$ the outward radial vector $R$ is tangent to the $\xi$
directions, so $S^*(M\times\Sphere)$ is oriented there by
$\iota_R(\omega_M^n/n!)\wedge d\tau\wedge d\theta=\sigma\wedge d\tau\wedge d\theta$, where
$\sigma$ is the Liouville orientation of $S^*M$. The outward normal of $S^*_+$ at
$\tau=0$ is $-\pa_\tau$, and $\iota_{-\pa_\tau}(\sigma\wedge d\tau\wedge d\theta)=
(-1)^{2n}\sigma\wedge d\theta$ since $\sigma$ has degree $2n-1$.
\end{proof}

We use the Atiyah--Singer theorem in the following form, with $\Td(M)$ a closed form
representing the Todd class of $TM\otimes\C$ (by Stokes' theorem the right side does
not depend on the choice of form).

\begin{theorem}[Atiyah--Singer, cosphere form]\label{eta:thm:AS}
If $W$ is a compact manifold without boundary and $P\in\Psido{0}(W;\C^N)$ is elliptic
with principal symbol $p\in\CI(S^*W;\GL(N,\C))$ then
\[
\ind(P)=\int_{S^*W}\Td(W)\wedge\Chodd(p),
\]
where $S^*W$ carries the Liouville orientation and only the component of degree $2\dim W-1$
is integrated.
\end{theorem}

This is the case of a point base of the families index theorem of \cref{ch:families},
once the relative Chern character of the symbol class in $K_c(T^*W)$ is transgressed to
the cosphere bundle; the normalization of $\Chodd$ is the one forced by Bott
periodicity, \cref{eta:prop:toeplitz-chern}. For $W=\Sphere$ it is the elementary formula
$\ind P=-\wind\det p(\cdot,1)+\wind\det p(\cdot,-1)$, which follows from
\eqref{eta:eq:toeplitz-index} by splitting $P$ into its positive and negative Toeplitz
parts (see \cref{ch:bott}); with the Liouville orientations ($d\theta$ on $\xi=1$,
$-d\theta$ on $\xi=-1$) this agrees with the theorem, which checks the orientation
convention.

\section{The index formula}\label{eta:sec:index}

For a fully elliptic $A\in\PsiptT{0,0}(M\times\Sphere;\C^N)$ with symbols
$(a,\beta)$, $b=\sigma_0(\beta)$, set
\begin{equation}\label{eta:eq:I}
I(A)=\int_{S^*M}\Td(M)\wedge\bigl(\eta(\beta)+\gamma(b)\bigr)
+\int_{S^*_+}\Td(M)\wedge\Chodd(a),
\end{equation}
where $\Td(M)$ is pulled back to $S^*M$ and to $S^*_+$, the orientations are those of the
previous section, and only components of top degree are integrated. Note that $I(A)$
depends only on the joint symbol of $A$.

\begin{theorem}\label{eta:thm:index}
For every fully elliptic $A\in\PsiptT{0,0}(M\times\Sphere;\C^N)$,
\[
\ind(A)=\int_{S^*M}\Td(M)\wedge\bigl(\eta(\beta)+\gamma(b)\bigr)
+\int_{S^*_+(M\times\Sphere)}\Td(M)\wedge\Chodd(a).
\]
\end{theorem}

For $M$ a point (formally $n=0$) the formula reduces to \eqref{eta:eq:toeplitz-index},
$S^*_+$ being the circle oriented by $d\theta$. The proof occupies the rest of this
section. We write $\Delta(A)=\ind(A)-I(A)$.

\begin{proposition}\label{eta:prop:homotopy}
If $(a_t,\beta_t)$, $t\in[0,1]$, is a smooth family of compatible pairs of fully
elliptic symbols then $I$ is constant in $t$. Consequently $\Delta(A)$ depends only on
the homotopy class of the joint symbol of $A$ through fully elliptic symbols.
\end{proposition}

\begin{proof} Write $T=\Td(M)$, a closed even form, $X=S^*M$, and
$I(t)=I_1(t)+I_2(t)$ for the two terms. By \cref{eta:lem:transgression},
$\pa_t\Chodd(a_t)=d\Theta(a_t)$, so by Stokes' theorem and \cref{eta:lem:orientation}
\[
I_2'(t)=\int_{S^*_+}T\wedge d\Theta(a_t)=\int_{X\times\Sphere}T\wedge\Theta(b_t)
=\int_XT\wedge\int_\Sphere\Theta(b_t),
\]
where $\Theta(b_t)$ is the restriction of $\Theta(a_t)$ to the boundary, given by the same
formula in terms of $b_t$. Now apply \cref{eta:thm:variation} to the family $\beta$
on $X\times[0,1]$, with total differential $D=d+dt\,\pa_t$. Writing
$\eta(\beta)+\gamma(b)=\zeta(t)+\zeta'(t)\wedge dt$ and, by
\cref{eta:lem:transgression} applied on $X\times\Sphere\times[0,1]$,
\[
\int_\Sphere\Chodd(b)=\int_\Sphere\Chodd(b_t)+\int_\Sphere\bigl(\Theta(b_t)\wedge dt\bigr)
=\int_\Sphere\Chodd(b_t)-\Bigl(\int_\Sphere\Theta(b_t)\Bigr)\wedge dt
\]
(moving $dt$ past $d\theta$), the $dt$-component of \eqref{eta:eq:variation-total} reads
$-\pa_t\zeta+d_X\zeta'=\int_\Sphere\Theta(b_t)$, using $dt\wedge\pa_t\zeta=-\pa_t\zeta\wedge dt$
($\zeta$ is odd). Since $T$ is closed and $X$ is compact without boundary,
\[
I_1'(t)=\int_XT\wedge\pa_t\zeta=-\int_XT\wedge\int_\Sphere\Theta(b_t)=-I_2'(t).
\]
The second statement follows from \ref{eta:P2} and \ref{eta:P4}: a homotopy of joint
symbols is realized by a smooth family of fully elliptic operators, along which the
index is constant.
\end{proof}

We use direct sums and stabilization freely: if $A$ acts on $\C^N$-valued functions,
$A\oplus\Pi$ acting on $\C^{N+K}$-valued functions has the same index and the same $I$,
and both $\ind$ and $I$ are additive under direct sums (all the forms involved are
traces of block-diagonal expressions).

\begin{lemma}\label{eta:lem:multiplicative}
If $A,A'\in\PsiptT{0,0}(M\times\Sphere;\C^N)$ are fully elliptic then
$\Delta(AA')=\Delta(A)+\Delta(A')$.
\end{lemma}

\begin{proof} The index is additive under composition of Fredholm operators. Let
$R_t$ be the constant unitary $\begin{psmallmatrix}\cos t&-\sin t\\ \sin t&\cos t\end{psmallmatrix}
\otimes\Id_N$ on $\C^{2N}$; it commutes with $\Pi$ and is an invertible element of the
algebra. Then
\[
C_t=R_t(A\oplus\Pi)R_t^{-1}(\Pi\oplus A'),\qquad 0\le t\le\pi/2,
\]
is a smooth family of fully elliptic operators with $C_0=A\oplus A'$ and
$C_{\pi/2}=\Pi\oplus AA'$. By \cref{eta:prop:homotopy} and additivity,
$I(A)+I(A')=I(\Pi\oplus AA')=I(AA')$, since $I(\Pi)=0$ ($\beta(\Pi)$ is constant and
$a=\Id$).
\end{proof}

\begin{lemma}\label{eta:lem:base-multiplication}
If $h\in\CI(M;\GL(N,\C))$ and $H$ denotes multiplication by $h$ on
$\Hardy(M\times\Sphere;\C^N)$, then $H$ is fully elliptic and $\ind(H)=I(H)=0$.
\end{lemma}

\begin{proof} $H$ is invertible, so $\ind(H)=0$. Its symbols are $a=h$ (pulled back to
$S^*_+$) and $\beta=\Pi h\Pi$, constant in $\xi$ and $\theta$. Hence $\gamma(b)=0$ and, by
\cref{eta:lem:circle}(iii), $\eta(\beta)=\frac12\Chodd(h)$ pulled back from $M$. The form
$\Td(M)\wedge\Chodd(h)$ is pulled back from $M$, which has dimension $n$, whereas
$\dim S^*_+=2n+1>n$ and $\dim S^*M=2n-1$; so both integrals in \eqref{eta:eq:I} vanish
when $n\ge2$. If $n=1$ the second still vanishes, and the first vanishes because
$S^*M\cong M\times\{\pm1\}$ with opposite orientations on the two components.
\end{proof}

\begin{lemma}\label{eta:lem:compression}
Let $P\in\Psido{0}(M\times\Sphere;\C^N)$ be elliptic with principal symbol $p$ equal to
$\Id$ outside a compact subset of $\{\tau>0\}\subset S^*(M\times\Sphere)$. Then
$\Pi P\Pi\in\PsiptT{0,0}$ is fully elliptic with $a=p|_{S^*_+}$, $\beta=\Pi$, and
\[
\ind(\Pi P\Pi)=\ind(P)=I(\Pi P\Pi).
\]
\end{lemma}

\begin{proof} By \ref{eta:P5}, $\beta(P)$ is multiplication by $p|_{\tau=0}=\Id$, so
$\beta(\Pi P\Pi)=\Pi$ and $\sigma(\Pi P\Pi)=p$ on $S^*_+$. With respect to
$L^2=\Hardy\oplus\Hardy^\perp$ write $P=\Pi P\Pi+(\Id-\Pi)P(\Id-\Pi)+K$,
$K=\Pi P(\Id-\Pi)+(\Id-\Pi)P\Pi$. Both symbols of $K$ vanish ($\sigma(\Pi)\sigma(\Id-\Pi)=0$
and $\Pi\,\Id\,(\Id-\Pi)=0$), and both symbols of $(\Id-\Pi)P(\Id-\Pi)-(\Id-\Pi)$ vanish
(on $S^*_-$, $p=\Id$). By \ref{eta:P3} these are compact, so $\ind(P)=\ind(\Pi P\Pi)
+\ind\bigl((\Id-\Pi)+\text{compact}\bigr)=\ind(\Pi P\Pi)$. On the other hand
$\eta(\Pi)=0$ and $\gamma(\Id)=0$, so $I(\Pi P\Pi)=\int_{S^*_+}\Td(M)\wedge\Chodd(p)$.
Since $p=\Id$ on $S^*_-$, $\Chodd(p)=0$ there, and $\Td(M\times\Sphere)=\Td(M)$ (as
$T(M\times\Sphere)\otimes\C=\pi^*(TM\otimes\C)\oplus\C$); so \cref{eta:thm:AS} on
$M\times\Sphere$ gives $\ind(P)=I(\Pi P\Pi)$.
\end{proof}

We also need the weak contractibility of the normalized Toeplitz group from
\cref{ch:bott}, in the following form. If $\beta'\in\CI(X;\GT(\Sphere;\C^N))$, $X$
compact, satisfies the normalization $\sigma_0(\beta')(x,1)=\Id$ (at $\theta=0$), then
after stabilization $\C^N\subset\C^K$, $\beta'\oplus\Pi$ is smoothly homotopic through
$\CI(X;\GT(\Sphere;\C^K))$ to the constant family $\Pi$. Indeed \cref{ch:bott} gives
such a homotopy in the group stabilized by the smoothing operators on an auxiliary
manifold $Y$, with $\C^N$ embedded as the span of the first $N$ eigenfunctions of a
Laplacian on $Y$; compressing the homotopy by the spectral projection onto the first
$K$ eigenfunctions gives, for $K$ large, a homotopy of the required kind, since the
compressions converge in the $\CI$ topology and invertibility is an open condition,
uniformly on the compact parameter space.

\begin{proof}[Proof of \cref{eta:thm:index}] Let $A$ be fully elliptic with symbols
$(a,\beta)$, $X=S^*M$. We show $\Delta(A)=0$ in four steps; each step replaces $A$ by
an operator with the same $\Delta$.

\emph{Homotopies with prescribed boundary values.} Choose a collar
$\pa S^*_+\times[0,\epsilon)_r$ of the boundary of $S^*_+$ and $\phi\in\CI([0,\epsilon))$ with
$\phi=0$ near $0$ and $\phi=1$ on $[\epsilon/2,\epsilon)$. If $b_t$ is a smooth homotopy of
$b_0=a|_{\pa S^*_+}$ through maps into $\GL(N,\C)$, then
\[
a_t(q,r)=b_{t(1-\phi(r))}(q)\,b_0(q)^{-1}a(q,r)\ \text{in the collar},\qquad
a_t=a\ \text{elsewhere},
\]
is a smooth homotopy of $a$ through invertible symbols with $a_t|_{\pa S^*_+}=b_t$. So any
smooth homotopy $\beta_t$ of $\beta$ through $\CI(X;\GT)$ extends to a homotopy of fully
elliptic compatible pairs $(a_t,\beta_t)$, which does not change $\Delta$
(\cref{eta:prop:homotopy}).

\emph{Step 1: $\beta$ constant in $\theta$.} Let $g(x)=b(x,1)\in\GL(N,\C)$ and
$\beta'=g^{-1}\beta$ (the matrix $g(x)^{-1}$ acting on $\Hardy\otimes\C^N$, which
commutes with $\Pi$). Then $\beta'$ is normalized, so after stabilization there is a
homotopy $\beta'_t$ from $\beta'$ to $\Pi$, and $\beta_t=g\beta'_t$ is a homotopy from $\beta$
to $\Pi g\Pi$. After extending it as above we may assume that $\beta(x)=\Pi g(x)\Pi$
with $g\in\CI(X;\GL(N,\C))$; then $b(x,\theta)=g(x)$.

\emph{Step 2: $\beta$ pulled back from $M$.} The submanifold $S^*_+\cap\{\theta=0\}$ is
a closed ball bundle over $M$ with boundary $X$, and $a$ restricted to it extends $g$.
Contracting the balls linearly to their centres (in the stereographic picture,
$v\mapsto(1-t)v$) gives a smooth homotopy $g_t$, $g_0=g$, from $g$ to $h\circ\pi_M$,
where $h(m)$ is the value of $a$ at the centre over $(m,0)$. The family
$\beta_t=\Pi g_t\Pi$ consists of invertible operators; extending as above we may assume
$\beta=\Pi h\Pi$ with $h\in\CI(M;\GL(N,\C))$.

\emph{Step 3: $\beta=\Pi$.} Let $H$ be multiplication by $h$ and $A'=H^{-1}A$. Then
$\beta(A')=\Pi$ and $a'=h^{-1}a$ satisfies $a'|_{\pa S^*_+}=\Id$. By
\cref{eta:lem:multiplicative,eta:lem:base-multiplication},
$\Delta(A)=\Delta(H)+\Delta(A')=\Delta(A')$.

\emph{Step 4: reduction to an ordinary operator.} Let $\kappa\in\CI([0,\epsilon))$ be
non-decreasing with $\kappa=0$ on $[0,\epsilon/2]$ and $\kappa(r)=r$ near $\epsilon$, and
let $\Phi_t(q,r)=(q,(1-t)r+t\kappa(r))$ in the collar, $\Phi_t=\Id$ elsewhere. Then
$a'\circ\Phi_t$ is a homotopy of invertible symbols with boundary value $\Id$, and
$a''=a'\circ\Phi_1$ equals $\Id$ near $\pa S^*_+$. Keeping $\beta=\Pi$, this is a homotopy
of fully elliptic pairs. Extend $a''$ by $\Id$ to a smooth function $p$ on
$S^*(M\times\Sphere)$ and choose $P\in\Psido{0}(M\times\Sphere;\C^N)$ with principal symbol
$p$. Then $\Pi P\Pi$ has the joint symbol $(a'',\Pi)$, so by
\cref{eta:lem:compression} and \cref{eta:prop:homotopy}, $\Delta(A')=\Delta(\Pi P\Pi)=0$.
\end{proof}

\begin{remark}\label{eta:rem:finite-rank}
The theorem contains the Atiyah--Singer formula on $M$. Let $\pi_{\le k}$ be the
projection onto the span of $e^{ij\theta}$, $0\le j\le k$, and let
$Q\in\Psido{0}(M;\C^N\otimes\C^{k+1})$ be elliptic, $\C^{k+1}\cong\Ran\pi_{\le k}$. The
operator $A=Q\pi_{\le k}+(\Pi-\pi_{\le k})$ is in $\PsiptT{0,0}$ (compare the compressions
in \cref{ch:product}) with $a=\Id$ and $\beta=\sigma(Q)\pi_{\le k}+\Pi-\pi_{\le k}$, a
finite-rank perturbation of $\Pi$; so $b=\Id$, $\gamma(b)=0$ and
$\eta(\beta)=\Chodd(\sigma(Q))$. Since $\ind(A)=\ind(Q)$, \cref{eta:thm:index} reduces to
\cref{eta:thm:AS} on $M$ with the Liouville orientation of $S^*M$, consistently with
\cref{eta:lem:orientation}. Conversely, one might try to prove the theorem by
deforming every $A$ to this form; but a fully elliptic $A$ need not be deformable to one
with $a=\Id$, because the obstruction, the index bundle over $M$ of the Toeplitz family
with symbol $a$ restricted to the centre section, need not vanish (it only vanishes
after pull-back to $S^*M$). This is why the proof above deforms in the opposite
direction and uses the Atiyah--Singer theorem on $M\times\Sphere$.
\end{remark}

\begin{remark}\label{eta:rem:gamma}
The symbolic term $\gamma(b)$ cannot be omitted. Take $n=1$, $M=\Sphere_x$, $N=1$, and
symbols $a_t=\exp(it\chi\sin(x+\theta))$ on $S^*_+$, where $\chi$ is a function on the
fibres of $S^*_+$ (half-circles) equal to $1$ at the end $\xi=+1$ and to $0$ at the end
$\xi=-1$. Thus $b_t=\exp(it\sin(x+\theta))$ on the component $\xi=+1$ of $S^*M$ and
$b_t=1$ on $\xi=-1$. Since $(a_t)$ starts at $a_0=1$, invertible base families $\beta_t$
with symbols $b_t$ and $\beta_0=\Pi$ exist (lift the homotopy as in \cref{ch:index}). In
degree one, with $\Td=1$ there,
\[
\int_{S^*M}\gamma(b_t)=c_0\int_0^{2\pi}\frac{i}{4\pi}\int_0^{2\pi}
(it\cos(x+\theta))^2\,d\theta\,dx=\frac{i}{2\pi}\cdot\frac{i}{4\pi}\cdot(-t^2)\cdot2\pi^2
=\frac{t^2}4,
\]
which depends on $t$. Since $I$ is constant along the family by
\cref{eta:prop:homotopy}, the remaining terms of $I$ vary by $-t^2/4$, so the formula
without $\gamma$ cannot hold for all $t$.
\end{remark}

\section{Non-Toeplitz product-type operators}\label{eta:sec:full}

The formula extends to the full product-type calculus, where both halves $S^*_\pm$ of the
blown-up cosphere bundle carry symbols. For a fully elliptic
$P\in\Psipt{0,0}(M\times\Sphere;\C^N)$ write $a_\pm=\sigma(P)|_{S^*_\pm}$,
$\beta=\beta(P)\in\CI(S^*M;\Psido{0}(\Sphere;\C^N))$ and $b_\pm=a_\pm|_{\pa S^*_\pm}$, regarded as
functions on $S^*M\times\Sphere$. The eta form $\eta(\beta)$ is defined exactly as before, with
$\Trbar$ on $\Psido{0}(\Sphere;\C^N)$. For a map $c\in\CI(S^*M\times\Sphere;\GL(N,\C))$ let
$\langle\Chodd(c)\rangle=\frac1{2\pi}\int_0^{2\pi}\Chodd(c(\cdot,\theta))\,d\theta$ be the
$\theta$-average of the odd Chern forms of the maps $c(\cdot,\theta)$ on $S^*M$.

\begin{proposition}\label{eta:prop:full}
For every fully elliptic $P\in\Psipt{0,0}(M\times\Sphere;\C^N)$,
\begin{multline*}
\ind(P)=\int_{S^*M}\Td(M)\wedge\bigl(\eta(\beta)+\langle\Chodd(b_-)\rangle+\gamma(b_+)-\gamma(b_-)\bigr)\\
+\int_{S^*_+}\Td(M)\wedge\Chodd(a_+)+\int_{S^*_-}\Td(M)\wedge\Chodd(a_-),
\end{multline*}
with $S^*_\pm$ oriented as subsets of $S^*(M\times\Sphere)$ and $S^*M$ Liouville
oriented.
\end{proposition}

\begin{proof} Let $J:\Hardy^\perp\longrightarrow\Hardy$ be the unitary
$(Ju)(\theta)=e^{-i\theta}u(-\theta)$, mapping $e^{-i(k+1)\theta}$ to $e^{ik\theta}$, and
\[
\mathcal U:L^2(M\times\Sphere;\C^N)\longrightarrow\Hardy(M\times\Sphere;\C^{2N}),\quad
\mathcal Uu=\bigl(\Pi u,J(\Id-\Pi)u\bigr).
\]
The reflection $\theta\mapsto-\theta$ is a fibre-preserving diffeomorphism, so it
preserves the product-type calculus, and the off-diagonal blocks $\Pi P(\Id-\Pi)$,
$(\Id-\Pi)P\Pi$ have vanishing principal symbol. Hence $P^\sharp=\mathcal UP\mathcal U^{-1}$
is a fully elliptic element of $\PsiptT{0,0}(M\times\Sphere;\C^{2N})$ with
\[
a^\sharp=a_+\oplus R^*a_-,\qquad\beta^\sharp=\mathcal U\beta\,\mathcal U^{-1},
\]
where $R:S^*_+\longrightarrow S^*_-$ is induced by $(\theta,\tau)\mapsto(-\theta,-\tau)$, and
$\ind(P)=\ind(P^\sharp)$. We evaluate $I(P^\sharp)$.

The map $R$ preserves the Liouville form $d\tau\wedge d\theta$ and the radial structure, so
it preserves orientations and $\int_{S^*_+}\Td\wedge\Chodd(R^*a_-)=\int_{S^*_-}\Td\wedge\Chodd(a_-)$.
The boundary symbol of $P^\sharp$ is $b_+\oplus b_-^R$ with
$b_-^R(x,\theta)=b_-(x,-\theta)$; the substitution $\theta\mapsto-\theta$ changes the sign of
$\mu$ and not of $\nu$, so $\gamma(b_-^R)=-\gamma(b_-)$. Finally
$\Trbar(\mathcal UQ\,\mathcal U^{-1})=\Trbar_{E_1}(Q)$ where $E_1=\mathcal U^{-1}E\,\mathcal U$
multiplies $e^{ij\theta}$ by $\lambda_j^s$ for $j\ge0$ and by $\lambda_{\abs j-1}^s$ for $j<0$.
This is an admissible family and, by \cref{eta:lem:reg-dependence},
$\Trbar_{E_1}(Q)=\Trbar(Q)+\TrR(QD)$ where $D=\frac{d}{ds}(E_1-E)|_{s=0}$ is the Fourier
multiplier by $\log\lambda_{\abs j-1}-\log\lambda_{\abs j}$ for $j<0$ and $0$ for
$j\ge0$, with symbol $\xi^{-1}$ (modulo lower order) in $\xi<0$ and rapidly
decreasing in $\xi>0$. With $Q=(\beta^{-1}d\beta)^{2k+1}$, \cref{eta:lem:circle}(i) gives
\[
\TrR(QD)=-\frac1{2\pi}\int_0^{2\pi}\tr\bigl(\nu_-^{2k+1}\cdot(-1)\bigr)d\theta
=\frac1{2\pi}\int_0^{2\pi}\tr\bigl((b_-^{-1}d_Xb_-)^{2k+1}\bigr)d\theta,
\]
so $\eta(\beta^\sharp)=\eta(\beta)+\langle\Chodd(b_-)\rangle$. Now apply
\cref{eta:thm:index} to $P^\sharp$.
\end{proof}

As a check, let $P\in\Psido{0}(M\times\Sphere;\C^N)$ be an ordinary elliptic operator.
Then $\beta$ is multiplication by $b=b_+=b_-$, so $\eta(\beta)=0$ by
\cref{eta:lem:circle}(iii) and the $\gamma$-terms cancel. For fixed $\theta$ the map
$b(\cdot,\theta)$ on $S^*M$ extends over the ball bundle $S^*_+\cap\{\theta\}$, hence is
homotopic to a map pulled back from $M$, so $\int_{S^*M}\Td\wedge\Chodd(b(\cdot,\theta))=0$
as in \cref{eta:lem:base-multiplication}. The formula therefore reduces to
\cref{eta:thm:AS} on $M\times\Sphere$. The asymmetric term $\langle\Chodd(b_-)\rangle$
reflects the fact that the regularization \eqref{eta:eq:circle-E} is not invariant under
$J$.

\section{Towards general product-type index theorems}\label{eta:sec:general}

In this final section we indicate how the results above should fit into a general
index theory for product-type families. Parts of it are conjectural, and are labelled
as such.

\subsection{Thom isomorphisms and the Todd class} For an oriented real vector bundle
$V\longrightarrow M$ of rank $r$ over a compact manifold, integration over the fibres gives
the Thom isomorphism $H^*_c(V)\longrightarrow H^{*-r}(M)$. For a complex vector bundle
there is also a Thom isomorphism $K_c(V)\longrightarrow K(M)$ in $K$-theory, extending Bott
periodicity (the case of trivial $V$). Both sides carry Chern characters, giving a square
\[
\begin{tikzcd}
K_c(V)\arrow[r,"\text{Thom}"]\arrow[d,"\Ch"'] & K(M)\arrow[d,"\Ch"]\\
H^{\mathrm{even}}_c(V)\arrow[r,"\text{Thom}"'] & H^{\mathrm{even}}(M)
\end{tikzcd}
\]
which does \emph{not} commute: the two composites differ by multiplication by a
characteristic class of $V$, which (with appropriate conventions) is the Todd class. This
is the origin of the factor $\Td(M)=\Td(TM\otimes\C)$ in the index formula. A direct
computation of the Chern character of the index bundle in the families index theorem
of \cref{ch:families} produces the Todd class in the same way.

\subsection{Product-type $K$-theory} Recall that for a fibration $\phi:M\longrightarrow B$ of
compact manifolds the families index theorem is the equality of two maps
\[
\ind_a=\ind_t:K^0_c(T^*(M/B))\longrightarrow K^0(B),
\]
the analytic index being defined by representing a class by the symbol of an elliptic
family $A\in\Psido{0}(M/B;\mathbb E)$ and taking the stabilized index bundle, the
topological index by embedding the fibration in a product $\Sphere^N\times B$.

Consider now a compound fibration of compact manifolds
\[
M\xrightarrow{\ \psi\ }Y\xrightarrow{\ \Phi\ }B,\qquad\phi=\Phi\circ\psi,
\]
and the families $A\in\Psipt{0,0}(M/B;\mathbb E)$ of operators on the fibres of $\phi$
which are of product type with respect to the fibration (over $\Phi^{-1}(b)$) induced by
$\psi$. Such a family has two symbols,
\[
\sigma(A)\in\CI\bigl([S^*(M/B),\psi^*S^*(Y/B)];\hom(\mathbb E)\bigr),\qquad
\beta(A)\in\Psido{0}\bigl(\pi^*M/S^*(Y/B);\mathbb E\bigr),
\]
the second being a family of operators on the fibres of $\psi$ parametrized by
$S^*(Y/B)$, and full ellipticity means invertibility of both. Then the analytic index
$\ind_a(A)\in K^0(B)$ is well defined by stabilization and depends only on
$(\sigma(A),\beta(A))$. Considering all compatible invertible pairs modulo bundle
isomorphism, stabilization and homotopy one expects an abelian group
$K^0_{c,\psi\text{-pt}}(T^*(M/B))$ (and an odd version $K^{-1}_{c,\psi\text{-pt}}$) on which
the analytic index is defined. Forgetting $\beta$ gives a map to $K^0_c(T^*(M/Y))$
(evaluation of $\sigma$ in the vertical directions of $\psi$), and the families index
of $\psi$, pulled back to $S^*(Y/B)$, gives a map
$K^0_c(T^*(M/Y))\longrightarrow K^0(S^*(Y/B))$ which vanishes on the image, expressing the
invertibility of $\beta$. Identifying $K^{-1}(S^*(Y/B))$ with homotopy classes of maps into
$G^{-\infty}$, i.e.\ with invertible base families with symbol the identity, one is led to:

\begin{remark}[Conjectural six-term sequence]\label{eta:rem:six-term}
There should be an exact sequence
\[
\begin{tikzcd}[column sep=small]
K^{-1}(S^*(Y/B))\arrow[r] & K^0_{c,\psi\text{-pt}}(T^*(M/B))\arrow[r,"\sigma"] &
K^0_c(T^*(M/Y))\arrow[d,"\ind"]\\
K^{-1}_c(T^*(M/Y))\arrow[u,"\ind"] & K^{-1}_{c,\psi\text{-pt}}(T^*(M/B))\arrow[l,"\sigma"'] &
K^0(S^*(Y/B))\arrow[l]
\end{tikzcd}
\]
and the analytic index should extend to an equality
$\ind_a=\ind_t:K^0_{c,\psi\text{-pt}}(T^*(M/B))\longrightarrow K^0(B)$ with a topological
index defined by embedding. A natural problem is to find a more traditional
realization of the group $K^0_{c,\psi\text{-pt}}(T^*(M/B))$, for which the six-term
sequence should serve as a guide. Groups of this type were subsequently constructed
by Melrose and Rochon; see \cite{Melrose-Rochon} for the closely related families index
for manifolds with boundary.
\end{remark}

\subsection{Product-type cohomology} The cohomological form of the index theorem
expresses the Chern character of the index in terms of a Chern character of the symbol.
\Cref{eta:thm:index} suggests what the latter should be in the product-type case: a
pair consisting of the Chern form of $\sigma$ and the eta form of $\beta$, related by
\cref{eta:thm:variation}. We formulate the corresponding cohomology in an abstract
setting (correcting the degree conventions of the original notes).

Let $Z$ be a compact manifold with boundary and $\psi:\pa Z\longrightarrow Y$ a fibration
over a compact manifold $Y$ without boundary, with compact oriented fibres of
dimension $f$, and let $\psi_*$ denote integration over the fibres with the `fibre
last' convention, so that $\psi_*d=d\psi_*$ and $\int_{\pa Z}=\int_Y\psi_*$ when $\pa Z$
is oriented by $Y$ followed by the fibres. Let $\iota:\pa Z\hookrightarrow Z$. Set
\[
C^k(Z,\psi)=\CI(Z;\Lambda^k)\oplus\CI(Y;\Lambda^{k-f-1}),\qquad
\mathcal D(u,\tau)=(du,\ \psi_*\iota^*u-d\tau).
\]
Then $\mathcal D^2=0$; a pair is closed if $du=0$ and $d\tau=\psi_*\iota^*u$, and exact if
it is of the form $(dv,\psi_*\iota^*v-d\varsigma)$. Denote the cohomology by $H^k(Z,\psi)$.

\begin{proposition}\label{eta:prop:psi-cohomology}
\begin{enumerate}[label=(\roman*)]
\item There is a long exact sequence
\begin{multline*}
\cdots\longrightarrow H^{k-f-1}(Y)\longrightarrow H^k(Z,\psi)\longrightarrow H^k(Z)\\
\xrightarrow{\ \psi_*\iota^*\ }H^{k-f}(Y)\longrightarrow H^{k+1}(Z,\psi)\longrightarrow\cdots.
\end{multline*}
\item If $\dim Z=n$, the pairing $\langle(u,\tau),[Z,\psi]\rangle=\int_Zu-\int_Y\tau$ on
$C^n(Z,\psi)$ vanishes on exact pairs, so defines a linear form on $H^n(Z,\psi)$.
\end{enumerate}
\end{proposition}

\begin{proof} (i) The sequence of complexes $0\longrightarrow\CI(Y;\Lambda^{\bullet-f-1})
\longrightarrow C^\bullet(Z,\psi)\longrightarrow\CI(Z;\Lambda^\bullet)\longrightarrow0$,
$\tau\mapsto(0,\tau)$ and $(u,\tau)\mapsto u$, is exact and compatible with the
differentials (up to the sign of $d$ on the first complex, which does not affect
cohomology). The connecting map sends $[u]$ to the class of $\psi_*\iota^*u$, since
$\mathcal D(u,0)=(0,\psi_*\iota^*u)$. (ii) $\int_Zdv-\int_Y(\psi_*\iota^*v-d\varsigma)=
\int_{\pa Z}v-\int_{\pa Z}v=0$ by Stokes' theorem.
\end{proof}

When $\psi$ is the identity of $\pa Z$ ($f=0$), $C^\bullet(Z,\psi)$ is the mapping cone of
$\iota^*$, and $H^k(Z,\psi)$ is the relative cohomology $H^k(Z,\pa Z)$, i.e.\ the
compactly supported cohomology of the interior (\cref{eta:ex:relative}).

In the setting of \cref{eta:thm:index}, with $Z=S^*_+$, $Y=S^*M$ and $\psi$ the projection
$S^*M\times\Sphere\longrightarrow S^*M$ ($f=1$), \cref{eta:thm:variation} says precisely that
\[
\Ch_{\mathrm{pt}}(A)=\bigl(\Td(M)\wedge\Chodd(a),\ -\Td(M)\wedge(\eta(\beta)+\gamma(b))\bigr)
\]
is closed, and \cref{eta:thm:index} becomes
\[
\ind(A)=\bigl\langle\Ch_{\mathrm{pt}}(A),[S^*_+,\psi]\bigr\rangle.
\]
The proof of \cref{eta:prop:homotopy} amounts to showing that the class of
$\Ch_{\mathrm{pt}}(A)$ is homotopy invariant. In general one expects a Chern character
$K^0_{c,\psi\text{-pt}}(T^*(M/B))\longrightarrow H^{\mathrm{even}}_{c,\psi\text{-pt}}(T^*(M/B))$,
given by the pair (relative Chern form of $\sigma$, eta form of $\beta$) and defined
using connections, mapping the $K$-theory six-term sequence of \cref{eta:rem:six-term}
to the corresponding cohomological one, together with a push-forward
$H^{\mathrm{even}}_{c,\psi\text{-pt}}(T^*(M/B))\longrightarrow H^{\mathrm{even}}(B)$ when the
fibres are oriented, so that $\Ch(\ind_a(A))$ is the push-forward of $\Td\cdot\Ch_{\mathrm{pt}}(A)$.

\subsection{The determinant line bundle}\label{eta:sec:determinant} The numerical index
is the degree-zero part of the Chern character of the index bundle. In the odd case the
degree-one part may be thought of as the `spectral flow' of the phase of a determinant.
In the even case the degree-two part corresponds to the determinant line bundle, which we
now construct; the degree-three part in the odd case corresponds to the curvature of a
gerbe, and the degree-four part in the even case to a $2$-gerbe.

Let $\phi:M\longrightarrow B$ be a fibration with compact fibres $Z_b$ and let
$P\in\Psido{0}(M/B;E_+,E_-)$ be an elliptic family of numerical index $0$. Then each
$P_b$ can be perturbed by a smoothing operator to be invertible, so the sets
\[
\mathcal P_b=\bigl\{P_b+Q_b;\ Q_b\in\Psido{-\infty}(Z_b;E_+,E_-),\ (P_b+Q_b)^{-1}\in
\Psido{0}(Z_b;E_-,E_+)\bigr\}
\]
are non-empty. Each $\mathcal P_b$ is a principal homogeneous space for the group
$G^{-\infty}(Z_b;E_-)$ acting on the left: if $A,A'\in\mathcal P_b$ then
$A'=(\Id+R)A$ with $\Id+R=A'A^{-1}$, $R=(A'-A)A^{-1}$ smoothing. Thus
$\mathcal P=\bigcup_b\mathcal P_b\longrightarrow B$ behaves like a principal bundle, although
the structure groups $G^{-\infty}(Z_b;E_-)$ form a bundle of groups rather than a fixed
group. (Using exhausting families of smoothing projections, as in the stabilization of the
index bundle in \cref{ch:index}, one sees that the components of the space of sections of
this bundle of groups are canonically $K^{-1}(B)$, so it is `weakly trivial'.) The
Fredholm determinant is a character of each $G^{-\infty}(Z_b;E_-)$, so we may form the
associated line bundle.

\begin{proposition}\label{eta:prop:det-bundle}
Let $\operatorname{Det}_b=(\mathcal P_b\times\C)/\sim$, where
$(A,z)\sim\bigl((\Id+R)A,\det(\Id+R)^{-1}z\bigr)$ for $\Id+R\in G^{-\infty}(Z_b;E_-)$. Then
$\operatorname{Det}=\bigcup_b\operatorname{Det}_b$ is a smooth line bundle over $B$, and
\[
s(b)=\bigl[A,\det(A^{-1}P_b)\bigr],\qquad A\in\mathcal P_b,
\]
is a well-defined smooth section of $\operatorname{Det}$ which vanishes exactly at the points
$b$ where $P_b$ is not invertible.
\end{proposition}

\begin{proof} The relation is an equivalence relation because $\det$ is multiplicative.
For fixed $A_0\in\mathcal P_b$, every element of $\mathcal P_b$ is uniquely $(\Id+R)A_0$, so
$z\mapsto[A_0,z]$ is a bijection $\C\longrightarrow\operatorname{Det}_b$, and the linear structures
induced by different choices of $A_0$ agree. Local triviality: if $A_0=P_{b_0}+Q_0\in\mathcal
P_{b_0}$, choose a local trivialization of the fibration near $b_0$ and transport $Q_0$ to a
smooth family $Q_b$; since invertibility is open, $P_b+Q_b\in\mathcal P_b$ for $b$ near
$b_0$, giving a local smooth section of $\mathcal P$, hence a local trivialization, and the
transition functions are $\det$ of smooth families in $G^{-\infty}$, hence smooth. For
$s$: $A^{-1}P_b=\Id-A^{-1}Q_b\in\Id+\Psido{-\infty}$ has a Fredholm determinant, which is
smooth in $b$ for a smooth local section $A$. Replacing $A$ by $(\Id+R)A$ gives
$\det(A^{-1}(\Id+R)^{-1}P_b)=\det(\Id+R)^{-1}\det(A^{-1}P_b)$, using the multiplicativity
of $\det$ and its invariance under conjugation by $A$; so $s$ is well defined.
Finally $\det(A^{-1}P_b)\ne0$ if and only if $A^{-1}P_b$, equivalently $P_b$, is invertible.
\end{proof}

When $P_b$ is invertible, $s(b)=[P_b,1]$. The section $s$ is the determinant of the family.
The curvature of natural connections on $\operatorname{Det}$ (Quillen, Bismut--Freed)
represents the degree-two part of $-\Ch(\ind P)$ (compare
\cref{eta:rem:determinant-connection}); with the opposite convention for the
equivalence relation one obtains the dual line bundle.

\section{Exercises}\label{eta:sec:exercises}

\begin{exercise}\label{eta:ex:defect-shift}
Let $U=\Pi e^{i\theta}\Pi$ and $L=\Pi e^{-i\theta}\Pi$ be the shifts on $\Hardy$. Show
that $[L,U]$ is the projection onto the constants, and verify the trace-defect formula
$\Trbar([L,U])=\TrR((D_EL)U)=1$ directly from \cref{eta:lem:circle}.
\end{exercise}

\begin{exercise}\label{eta:ex:not-inner}
Show that $D_E$ is not an inner derivation of $\Psido{\Z}(\Sphere)$. (Hint: if
$D_E=[C,\cdot]$ with $C\in\Psido{k}(\Sphere)$, compare the principal symbols of
$D_Ef=-i\xi^{-1}f'+\dots$ and $[C,f]$, $f\in\CI(\Sphere)$; show that $k\le0$ is forced and
that $k=0$ would require $\pa_\xi\sigma_0(C)=\xi^{-1}$ with $\sigma_0(C)$ homogeneous of
degree $0$.)
\end{exercise}

\begin{exercise}\label{eta:ex:change-eta}
Let $E'$ be another admissible family on $\Sphere$ with $D=D(E,E')\in\Psido{-1}(\Sphere)$ of
principal symbol $\delta(\theta,\xi)$. Show that the eta forms defined using $E'$ are
\[
\eta'(B)=\eta(B)-\frac1{2\pi}\sum_kc_k\int_0^{2\pi}\tr(\nu^{2k+1})\,\delta(\theta,1)\,d\theta,
\]
so that they differ from $\eta$ by a term local in the symbol of $B$. Deduce that
\cref{eta:thm:index} holds with $\eta'$ in place of $\eta$ provided $\gamma$ is replaced by
$\gamma$ minus this local term.
\end{exercise}

\begin{exercise}\label{eta:ex:chodd-product}
Let $g,h\in\CI(W;\GL(N,\C))$. Using the homotopy
$R_t(g\oplus\Id)R_t^{-1}(\Id\oplus h)$ and \cref{eta:lem:transgression}, show that
$\Chodd(gh)-\Chodd(g)-\Chodd(h)$ is exact, and that its degree-one part vanishes.
\end{exercise}

\begin{exercise}\label{eta:ex:circle-AS}
Show that for elliptic $P\in\Psido{0}(\Sphere;\C^N)$,
$\ind(P)=\int_{S^*\Sphere}\Chodd(\sigma_0(P))$ with the Liouville orientation, by writing
$P$, up to compact operators, as $\Pi P_+\Pi+(\Id-\Pi)P_-(\Id-\Pi)$ and applying
\eqref{eta:eq:toeplitz-index} to both parts (for the negative part use $J$ from the proof of
\cref{eta:prop:full}).
\end{exercise}

\begin{exercise}\label{eta:ex:orientation}
Verify directly in local coordinates that for $n=1$ the Liouville orientation of
$S^*M\cong M\times\{\pm1\}$ is $\pm dx$, and that the orientation of $\pa S^*_+$ is
$\pm dx\wedge d\theta$ on the two components.
\end{exercise}

\begin{exercise}\label{eta:ex:relative}
Check that when $\psi:\pa Z\longrightarrow\pa Z$ is the identity, $H^k(Z,\psi)$ is
naturally isomorphic to $H^k(Z,\pa Z)=H^k_c(Z\setminus\pa Z)$. (Hint: given a closed pair
$(u,\tau)$, use a collar and a cut-off $\chi$ equal to $1$ near $\pa Z$ to replace it by the
cohomologous pair $(u-d(\chi\tau),0)$, with $u-d(\chi\tau)$ vanishing near $\pa Z$.)
\end{exercise}

\begin{exercise}\label{eta:ex:les}
Give the details of \cref{eta:prop:psi-cohomology}(i), and show that the connecting map
$H^k(Z)\longrightarrow H^{k-f}(Y)$ is induced by $\psi_*\iota^*$.
\end{exercise}

\begin{exercise}\label{eta:ex:det-section}
In the setting of \cref{eta:prop:det-bundle}, suppose $P_b=P_0+K_b$ where $P_0$ is a
fixed invertible family (in a local trivialization of the fibration) and
$K_b\in\Psido{-\infty}$. Show that in the trivialization $z\mapsto[P_0,z]$ of
$\operatorname{Det}$ the determinant section is $s(b)=\det(\Id+P_0^{-1}K_b)$, and that
$\operatorname{Det}$ is trivial if every $P_b$ is invertible.
\end{exercise}

\section*{Notes}

Regularized traces of this kind were studied by Kontsevich and Vishik, and the
trace-defect formula appears in the treatment of the b-trace in \cite{Melrose-APS}; the
residue trace is due to Wodzicki
and Guillemin \cite{Wodzicki-residue,Guillemin-residue}, and the meromorphic
continuation to Seeley \cite{Seeley}. In the APS index theorem the eta invariant arises in
the same way, as a regularized trace of the indicial family \cite{Melrose-APS}; the eta
forms here are the corresponding objects for the odd Chern character, analogous to the
eta forms of Bismut and Cheeger \cite{Bismut-Cheeger}, and appear in the families index
theorem of Melrose and Rochon \cite{Melrose-Rochon}.

The main departures from the original lecture notes are as follows. We have fixed all
normalizations: with $E(s)$ of order $+s$ the residue trace is minus the Wodzicki residue,
the trace-defect formula reads $\Trbar_E([A,B])=\TrR(BD_EA)$ with
$D_EA=\frac{d}{ds}[E(s),A]|_{s=0}$, and the constants of the odd Chern character are
determined by the Toeplitz index (\cref{eta:prop:toeplitz-chern}). The variation formula
for the eta forms, \cref{eta:thm:variation}, has been corrected; it contains a symbolic
term $\gamma$ which must also appear in the index formula (\cref{eta:rem:gamma}), and with
these normalizations the eta form enters with coefficient $+1$ rather than $-\frac12$.
The original proof of the index formula reduced to the Atiyah--Singer theorem on $M$ by
composing with the inverse of a bundle isomorphism; in the Toeplitz algebra such an
inverse need not exist (\cref{eta:rem:finite-rank}), and we have instead given a proof
based on the weak contractibility of the Toeplitz group and the Atiyah--Singer theorem on
$M\times\Sphere$. The non-Toeplitz case is deduced by a unitary doubling. In the product-type
cohomology the compatibility condition has been changed to $d\tau=\psi_*\iota^*u$ (push-forward
along the fibres), which is what the variation formula requires; the suggested
identification with the cohomology of the space obtained by collapsing the fibres of
$\psi$ refers to the pull-back condition and has been dropped. The original also asserts
that $D_E$ is the unique continuous outer derivation up to inner derivations and
constant multiples; we have not proved this. For the Todd class in the comparison of
Thom isomorphisms see Atiyah and Singer, \emph{The index of elliptic operators III}, and
for determinant line bundles Quillen and Bismut--Freed.
